% =====================================================================
%  广义对称性（Generalized Symmetry）课前预习讲义
%  课程：Generalized Symmetry
%  授课：Daniel Brennan
%  说明：本文件为自包含（self-contained）LaTeX 源文件。
%        请使用 XeLaTeX（推荐）或 LuaLaTeX 编译三次以生成正确的目录与交叉引用。
%        编译命令示例： xelatex 02-generalized-symmetry-daniel-brennan.tex
% =====================================================================
\documentclass[UTF8,a4paper,11pt]{ctexart}

\usepackage{amsmath}
\usepackage{amssymb}
\usepackage{amsthm}
\usepackage{geometry}
\usepackage{booktabs}
\usepackage{enumitem}
\usepackage{mathtools}
\usepackage{bm}
\usepackage{hyperref}

\geometry{left=2.4cm,right=2.4cm,top=2.6cm,bottom=2.6cm}

\hypersetup{
  colorlinks=true,
  linkcolor=blue,
  citecolor=blue,
  urlcolor=magenta,
  pdftitle={广义对称性（Generalized Symmetry）课前预习讲义},
  pdfauthor={课前自学材料},
  pdfsubject={Generalized Global Symmetries, Higher-Form Symmetry, Daniel Brennan}
}

% ---------------------------------------------------------------------
%  定理类环境（环境名用 ASCII，显示名用中文）
% ---------------------------------------------------------------------
\theoremstyle{plain}
\newtheorem{theorem}{定理}[section]
\newtheorem{proposition}{命题}[section]
\newtheorem{lemma}{引理}[section]

\theoremstyle{definition}
\newtheorem{definition}{定义}[section]
\newtheorem{example}{例}[section]
\newtheorem{exercise}{练习}[section]

\theoremstyle{remark}
\newtheorem{remark}{注}[section]
\newtheorem{keypoint}{要点}[section]
\newtheorem{warning}{注意}[section]

% ---------------------------------------------------------------------
%  记号宏
% ---------------------------------------------------------------------
\newcommand{\ext}{\mathrm{d}}          % 外微分 / 微分
\newcommand{\hs}{\star}                % Hodge 星号
\newcommand{\Tr}{\mathrm{Tr}}
\newcommand{\Pexp}{\mathbf{P}\exp}
\newcommand{\Link}{\mathrm{Link}}
\newcommand{\Vol}{\mathrm{vol}}
\newcommand{\ZZ}{\mathbb{Z}}
\newcommand{\RR}{\mathbb{R}}
\newcommand{\CC}{\mathbb{C}}
\newcommand{\Uone}{\mathrm{U}(1)}
\newcommand{\Ugrp}{\mathrm{U}}
\newcommand{\SUgrp}{\mathrm{SU}}
\newcommand{\SOgrp}{\mathrm{SO}}
\newcommand{\Op}{\mathcal{O}}
\newcommand{\Mfd}{\mathcal{M}}
\newcommand{\Lag}{\mathcal{L}}
\newcommand{\Ham}{\mathcal{H}}
\newcommand{\Ssurf}{\Sigma}
\newcommand{\PD}[2]{\frac{\partial #1}{\partial #2}}
\newcommand{\en}[1]{\textsf{#1}}       % 英文术语
\newcommand{\Uop}{U}
\newcommand{\Dop}{\mathcal{D}}
\newcommand{\pform}[1]{#1\text{-form}}

\allowdisplaybreaks[3]
\setlist{itemsep=2pt,parsep=2pt,topsep=3pt}

% =====================================================================
\title{\bfseries 广义对称性（Generalized Symmetry）课前预习讲义\\[8pt]
       \large 面向 Daniel Brennan 教授课程《Generalized Symmetry》}
\author{课前自学材料（简体中文，英文术语另列）}
\date{\today}

\begin{document}
\maketitle

\begin{abstract}
\noindent
本讲义为选修 Daniel Brennan 教授《广义对称性》（\en{Generalized Symmetry}）
课程的同学准备的\textbf{课前预习}材料。预设读者学过基础量子力学与量子场论
（正则量子化、路径积分、Feynman 规则），但\textbf{没有系统学习过广义对称性／
高阶形式对称性}（\en{higher-form symmetry}），也可能对微分形式语言不熟悉。

讲义围绕一条主线展开：\textbf{对称性 $=$ 拓扑算符}
（\en{symmetry} $=$ \en{topological operator}）。
这句话是整个领域的“重新表述”（\en{reformulation}），一旦接受它，
“为什么可以有 $p$-form 对称性”“为什么会有不可逆对称性”
就都变成自然的推论，而不是需要额外记忆的新概念。为此我们：
（一）先用最传统的 Noether 定理把守恒流推出来，
再\textbf{把它改写成拓扑算符}，让读者亲眼看到两种语言等价；
（二）\textbf{完整给出每一步数学推导}，包括微分形式、Stokes 定理、
Maxwell 理论中 1-form 对称性的全部细节，每一步都可用纸笔复核；
（三）强调物理图像与\textbf{可证伪的判据}
（“怎样判断某个算符是拓扑的”“怎样判断对称性被自发破缺”）。

所有英文术语首次出现时括号标注，并在第 \ref{sec:glossary} 节汇总为英汉术语表。
\end{abstract}

\vspace{0.4em}
\tableofcontents
\vspace{1.2em}

% =====================================================================
\section{导读}
\label{sec:guide}

\subsection{这门课到底在讲什么}

过去十年，量子场论（\en{quantum field theory}, QFT）中“对称性”的定义被
系统地扩充了。传统教科书里的对称性是：
\begin{center}
\emph{一个群 $G$ 作用在场上，使作用量不变；由 Noether 定理得到守恒流
$j^\mu$ 与守恒荷 $Q$。}
\end{center}
这个定义有三个隐含假设，而\textbf{每一个都可以去掉}：

\begin{enumerate}[label=\textbf{(A\arabic*)}, leftmargin=1.5cm]
  \item \textbf{荷作用在局域算符（点状对象）上}。
        去掉它 $\Rightarrow$ 荷可以作用在\textbf{线、面等拓展对象}上，
        得到 $p$-form 对称性（\en{$p$-form symmetry}）。
  \item \textbf{对称性作用是可逆的（群）}。
        去掉它 $\Rightarrow$ \textbf{不可逆对称性}
        （\en{non-invertible symmetry}），对称算符只满足融合代数
        （\en{fusion algebra}）而未必有逆。
  \item \textbf{对称性由连续变换的生成元定义}。
        去掉它 $\Rightarrow$ 离散、有限群甚至“范畴型”对称性
        （\en{categorical symmetry}）都被统一描述。
\end{enumerate}

把这三条统称为\textbf{广义（全局）对称性}
（\en{generalized global symmetry}）。它们的共同语言是：
\begin{equation}
  \boxed{\ \text{一个对称性} \;\Longleftrightarrow\;
  \text{一族\emph{拓扑}算符（在时空中支撑于某个子流形上）}\ }
  \label{eq:slogan}
\end{equation}

\subsection{为什么值得学}

这不只是形式上的重新包装，它带来了真正的新结果：

\begin{itemize}
  \item \textbf{禁闭的对称性判据}。4 维 $\SUgrp(N)$ Yang--Mills 理论有一个
        $\ZZ_N$ 的 1-form 中心对称性（\en{center symmetry}）；
        \textbf{禁闭（\en{confinement}）等价于这个 1-form 对称性未被自发破缺}，
        Wilson 圈的面积律就是它的序参量。这把“禁闭”从一个定性描述
        变成了一个对称性破缺的判断（第 \ref{sec:ssb} 节）。
  \item \textbf{新的 't Hooft 反常与相图约束}。高阶形式对称性之间的混合反常
        （\en{mixed anomaly}）给出对低能物理的严格限制，
        例如 4 维 Yang--Mills 在 $\theta$ 角处的相变结构。
  \item \textbf{“同一个局域理论”的全局结构分类}。$\Uone$ 与 $\Uone/\ZZ_N$
        规范理论有相同的局域动力学，却有不同的线算符谱；
        这种差别正是由 1-form 对称性的“规范化”（\en{gauging}）选择决定的。
  \item \textbf{不可逆对称性带来的选择定则}。例如 Kramers--Wannier 对偶
        在 2 维 Ising 临界点上就是一条不可逆拓扑线，
        它给出普通群论看不到的关联函数约束。
\end{itemize}

\subsection{本讲义的结构与建议读法}

本讲义有意把“旧语言”与“新语言”\textbf{并置}：

\begin{itemize}
  \item 第 \ref{sec:prelim}--\ref{sec:noether} 节完全是传统内容
        （群论、拉氏量、正则量子化、Noether 定理），
        但\textbf{结尾处做一次关键改写}（第 \ref{subsec:topological-rewrite} 小节），
        把 $Q$ 写成 $\int_\Ssurf \hs j$ 并证明它只依赖 $\Ssurf$ 的同调类。
        \textbf{这一步是整门课的枢纽}，请务必自己推一遍。
  \item 第 \ref{sec:operators}--\ref{sec:topological} 节介绍拓展算符与拓扑算符，
        把“对称性 $=$ 拓扑算符”这句话讲清楚。
  \item 第 \ref{sec:pform} 节做推广，第 \ref{sec:charges-defects} 节区分
        “荷”与“缺陷”。
  \item 第 \ref{sec:examples} 节是\textbf{全篇重心}：Ising 的 $\ZZ_2$、
        Maxwell 理论 1-form 对称性的完整推导、$\Uone$ 对偶直觉。
  \item 第 \ref{sec:ssb}--\ref{sec:further} 节是自发破缺、反常、术语表、
        困惑澄清、听课重点与文献路线。
\end{itemize}

\paragraph{三条阅读路径。}
\begin{description}[leftmargin=1.1cm,style=nextline]
\item[时间紧（约 2--3 小时）]
  第 \ref{subsec:forms} 小节（微分形式速成）$\to$
  第 \ref{subsec:topological-rewrite} 小节（$Q$ 的拓扑改写）$\to$
  第 \ref{sec:pform} 节的表 \ref{tab:pform}（$p$-form 对称性字典）$\to$
  第 \ref{subsec:maxwell} 小节（Maxwell 例子）$\to$
  第 \ref{sec:glossary} 节术语表 $\to$ 第 \ref{sec:focus} 节听课重点。
\item[时间充裕（约 8--10 小时）]
  按顺序全读并完成正文中的\textbf{练习}。重点动手项有四个：
  Noether 流的两个例子；Ward 恒等式；
  Maxwell 的 Gauss 律推导；
  Wilson 圈与 't Hooft 线的连接数计算。
\item[只想补数学语言]
  第 \ref{subsec:forms} 小节 $+$ 附录 \ref{app:forms}（微分形式与拓扑速查）。
\end{description}

\subsection{符号与约定}

\begin{itemize}
  \item 时空维数记作 $d$（含时间），度规号差取
        $\eta_{\mu\nu}=\mathrm{diag}(-,+,+,\dots,+)$，
        除特别说明外用自然单位 $\hbar=c=1$。
        \textbf{注意}：广义对称性文献中 $(-,+,\dots)$ 与 $(+,-,\dots)$
        两种号差都常见，欧氏与洛伦兹信号也混用；
        本讲义在需要时会明确指出。
  \item 时空流形记作 $\Mfd_d$（或 $X$）；其中的 $k$ 维闭子流形记作
        $\Ssurf_k,\ \Mfd_k,\ C$（曲线）等。
  \item 微分形式用 $\omega_{(p)}$ 表示 $p$-形式；外微分 $\ext$，
        Hodge 星号 $\hs$。
  \item “$q$-form 对称性”中的 $q$ 指\textbf{被作用的带荷对象的维数}
        （$q=0$ 即普通对称性作用在点状局域算符上）。
        这是 Gaiotto--Kapustin--Seiberg--Willett 的约定，全篇统一使用。
  \item 规范场 $A$ 视作 1-形式（$A=A_\mu\ext x^\mu$），
        场强 $F=\ext A$。规范群为 $\Uone$ 时 $A$ 是实的；
        非阿贝尔时 $A$ 取值于李代数。
\end{itemize}

\begin{warning}
本领域\textbf{约定极不统一}：$2\pi$ 因子的位置、
电磁对偶中 $e^2$ 的放法、“$p$-form”里 $p$ 指荷还是指算符维数、
反常的符号……都因作者而异。本讲义的策略是：
\textbf{每个关键公式都给出可自查的推导}，
这样即使老师用了别的约定，你也能当场把因子对上。
\end{warning}

% =====================================================================
\section{预备知识}
\label{sec:prelim}

本节复习后面要用的四样工具：群论语言、拉氏量形式、正则量子化，
以及\textbf{微分形式}。前三样多数读者已熟悉，可快速浏览；
第 \ref{subsec:forms} 小节即使学过也建议细读，
因为广义对称性的全部陈述都写在这套语言里。

\subsection{群论基础}
\label{subsec:group}

\begin{definition}[群]
集合 $G$ 配一个乘法 $G\times G\to G$，满足结合律、存在单位元 $e$、
每个元素 $g$ 存在逆元 $g^{-1}$。若还满足 $g_1g_2=g_2g_1$，称
\textbf{阿贝尔群}（\en{abelian group}）。
\end{definition}

后面反复出现的具体群：

\begin{itemize}
  \item $\ZZ_N=\{0,1,\dots,N-1\}$，加法模 $N$；等价地
        $\{1,\omega,\dots,\omega^{N-1}\}$，$\omega=e^{2\pi i/N}$，乘法。
        阿贝尔、有限、\textbf{离散}。
  \item $\Uone=\{e^{i\alpha}:\alpha\in[0,2\pi)\}$：阿贝尔、\textbf{连续}、紧致。
        它是本课程最重要的群。
  \item $\SUgrp(N)$：非阿贝尔、连续、紧致。其\textbf{中心}
        （\en{center}）为
        \begin{equation}
          Z\big(\SUgrp(N)\big)
          =\Big\{\omega\,\mathbf 1_N:\ \omega^N=1\Big\}\cong\ZZ_N .
          \label{eq:center}
        \end{equation}
        这个 $\ZZ_N$ 将在第 \ref{subsec:center} 小节变成 1-form 对称性。
\end{itemize}

\begin{definition}[表示]
群 $G$ 的（有限维）\textbf{表示}（\en{representation}）是同态
$\rho:G\to GL(V)$，即 $\rho(g_1g_2)=\rho(g_1)\rho(g_2)$。
物理上：场按某个表示变换，$\phi\to\rho(g)\phi$。
\end{definition}

\begin{example}[$\ZZ_N$ 的一维表示与“$N$ 度”]
\label{ex:Nality}
$\ZZ_N$ 的所有一维表示为 $\rho_n(\omega)=\omega^n$，$n=0,\dots,N-1$。
对 $\SUgrp(N)$ 的表示 $R$，其中心 $\omega\mathbf 1$ 作用为
$\rho_R(\omega\mathbf 1)=\omega^{n_R}\mathbf 1$，
$n_R\in\ZZ_N$ 称为该表示的 \textbf{$N$ 度}（\en{$N$-ality}）。
基本表示 $n=1$，伴随表示 $n=0$。
第 \ref{subsec:center} 小节会看到：$n_R$ 就是 Wilson 线在
1-form 中心对称性下的荷。
\end{example}

\paragraph{李代数与生成元。}对连续群，取 $g=e^{i\alpha^aT^a}$，
$T^a$ 为生成元，满足 $[T^a,T^b]=if^{abc}T^c$。
无穷小变换 $\delta\phi=i\alpha^aT^a\phi$。
对 $\Uone$ 只有一个生成元，$\delta\phi=i\alpha\phi$。

\paragraph{正规子群与商群。}若 $H\triangleleft G$，可作商 $G/H$。
这在“规范化对称性”（\en{gauging}）时至关重要：
规范化 $H$ 得到的理论，其对称性与 $G/H$ 有关，
且会\textbf{产生新的对偶对称性}（第 \ref{subsec:gauging} 小节）。

\begin{example}[$\Uone$ 与 $\Uone/\ZZ_N$]
$\ZZ_N\subset\Uone$（$N$ 次单位根），
$\Uone/\ZZ_N\cong\Uone$ 作为抽象群同构，
但作为\textbf{规范群}它们给出不同理论：
允许的电荷格（\en{charge lattice}）不同。
第 \ref{subsec:duality} 小节会回到这一点。
\end{example}

\subsection{拉氏量形式}
\label{subsec:lagrangian}

场论由作用量给出。设场为 $\phi^a(x)$（$a$ 标记场的种类与分量），
\begin{equation}
  S[\phi]=\int_{\Mfd_d}\ext^dx\;\Lag\big(\phi^a,\partial_\mu\phi^a\big).
  \label{eq:action}
\end{equation}

\begin{proposition}[Euler--Lagrange 方程]
$S$ 在 $\phi\to\phi+\delta\phi$（$\delta\phi$ 在边界上为零）下取极值的条件是
\begin{equation}
  \boxed{\ \PD{\Lag}{\phi^a}
  -\partial_\mu\!\left(\PD{\Lag}{(\partial_\mu\phi^a)}\right)=0\ }
  \label{eq:EL}
\end{equation}
\end{proposition}

\begin{proof}
\begin{align}
  \delta S
  &=\int\ext^dx\left[\PD{\Lag}{\phi^a}\delta\phi^a
    +\PD{\Lag}{(\partial_\mu\phi^a)}\,\partial_\mu\delta\phi^a\right]
  \notag\\
  &=\int\ext^dx\left[\PD{\Lag}{\phi^a}
    -\partial_\mu\PD{\Lag}{(\partial_\mu\phi^a)}\right]\delta\phi^a
    +\underbrace{\int\ext^dx\;\partial_\mu\!\left[
    \PD{\Lag}{(\partial_\mu\phi^a)}\delta\phi^a\right]}_{\text{边界项}=0} .
\end{align}
因 $\delta\phi^a$ 任意，方括号必须为零。
\end{proof}

\begin{example}[本讲义要用的三个拉氏量]
\label{ex:threelag}
\begin{enumerate}[label=(\arabic*)]
  \item \textbf{实标量}：$\Lag=-\tfrac12\partial_\mu\varphi\partial^\mu\varphi
        -V(\varphi)$。若 $V(\varphi)=V(-\varphi)$，有 $\ZZ_2$ 对称性
        $\varphi\to-\varphi$。
  \item \textbf{复标量}：$\Lag=-\partial_\mu\phi^*\partial^\mu\phi-V(|\phi|^2)$，
        有 $\Uone$ 对称性 $\phi\to e^{i\alpha}\phi$。
  \item \textbf{Maxwell}：$\Lag=-\dfrac{1}{4e^2}F_{\mu\nu}F^{\mu\nu}$，
        $F_{\mu\nu}=\partial_\mu A_\nu-\partial_\nu A_\mu$。
        \textbf{注意这里没有带电物质场}——这正是第
        \ref{subsec:maxwell} 小节能找到 1-form 对称性的原因。
\end{enumerate}
（号差 $(-,+,+,+)$；不同号差下整体符号会变，物理结论不变。）
\end{example}

\subsection{正则量子化简要}
\label{subsec:canonical}

正则量子化的三步：

\begin{enumerate}[label=\textbf{第 \arabic* 步}, leftmargin=2.1cm]
  \item \textbf{定义共轭动量}
        \begin{equation}
          \pi_a(x)=\PD{\Lag}{\dot\phi^a(x)} .
        \end{equation}
  \item \textbf{等时正则关系}（把 Poisson 括号换成交换子）
        \begin{equation}
          \big[\phi^a(t,\bm x),\ \pi_b(t,\bm y)\big]
          =i\,\delta^a_{\ b}\,\delta^{d-1}(\bm x-\bm y),
          \qquad
          [\phi,\phi]=[\pi,\pi]=0 .
          \label{eq:canonical}
        \end{equation}
  \item \textbf{Hamilton 量生成时间演化}
        \begin{equation}
          H=\int\ext^{d-1}x\,\big(\pi_a\dot\phi^a-\Lag\big),
          \qquad
          \dot{\Op}=i[H,\Op] .
        \end{equation}
\end{enumerate}

\begin{keypoint}
本讲义只用正则量子化做一件事：
\textbf{验证守恒荷 $Q$ 确实生成对称变换}（第 \ref{subsec:generates} 小节）。
这条验证是理解 \eqref{eq:slogan} 的必要一步：
“对称算符”$U=e^{i\alpha Q}$ 之所以是算符，
正因为 $Q$ 由场算符构成并作用在 Hilbert 空间上。
\end{keypoint}

\paragraph{路径积分表述。}后面讨论关联函数与 Ward 恒等式时会用
\begin{equation}
  \big\langle \Op_1(x_1)\cdots\Op_n(x_n)\big\rangle
  =\frac{1}{Z}\int\!\Dop\phi\;
  \Op_1(x_1)\cdots\Op_n(x_n)\;e^{iS[\phi]},
  \qquad
  Z=\int\!\Dop\phi\;e^{iS[\phi]} .
  \label{eq:pathintegral}
\end{equation}
在欧氏信号下 $e^{iS}\to e^{-S_E}$。
\textbf{广义对称性的大部分论证在欧氏信号、路径积分语言下最自然}，
因为“把算符支撑在任意子流形上并连续形变”在那里是显然的操作。

\subsection{微分形式速成}
\label{subsec:forms}

这是本节最重要的部分。广义对称性的所有陈述——
守恒律、算符支撑、连接数——都用微分形式写才简洁。

\subsubsection{\texorpdfstring{$p$}{p}-形式与楔积}

\begin{definition}[$p$-形式]
$d$ 维流形上的 \textbf{$p$-形式}（\en{$p$-form}）是完全反对称张量场，
写作
\begin{equation}
  \omega=\frac{1}{p!}\,\omega_{\mu_1\cdots\mu_p}\,
  \ext x^{\mu_1}\wedge\cdots\wedge\ext x^{\mu_p},
  \qquad
  \omega_{\cdots\mu_i\cdots\mu_j\cdots}=-\omega_{\cdots\mu_j\cdots\mu_i\cdots} .
\end{equation}
$p=0$ 即标量函数。因反对称性，$d$ 维中 $p>d$ 的形式恒为零。
\end{definition}

\textbf{楔积}（\en{wedge product}）满足
\begin{equation}
  \ext x^\mu\wedge \ext x^\nu=-\,\ext x^\nu\wedge \ext x^\mu,
  \qquad
  \omega_{(p)}\wedge\eta_{(q)}=(-1)^{pq}\,\eta_{(q)}\wedge\omega_{(p)} .
  \label{eq:wedge}
\end{equation}
特别地，奇数次形式与自身的楔积为零：$\omega_{(1)}\wedge\omega_{(1)}=0$。

\subsubsection{外微分}

\begin{definition}[外微分]
$\ext:\Omega^p\to\Omega^{p+1}$，
\begin{equation}
  \ext\omega=\frac{1}{p!}\,\partial_\nu\omega_{\mu_1\cdots\mu_p}\,
  \ext x^\nu\wedge \ext x^{\mu_1}\wedge\cdots\wedge \ext x^{\mu_p} .
\end{equation}
\end{definition}

两条核心性质：
\begin{equation}
  \boxed{\ \ext^2=0\ },
  \qquad
  \ext\big(\omega_{(p)}\wedge\eta_{(q)}\big)
  =\ext\omega_{(p)}\wedge\eta_{(q)}
  +(-1)^p\,\omega_{(p)}\wedge \ext\eta_{(q)} .
  \label{eq:dsquared}
\end{equation}
$\ext^2=0$ 的理由：$\ext^2\omega$ 的系数含
$\partial_\rho\partial_\nu\omega_{\cdots}$，
对 $\rho\nu$ 对称，而 $\ext x^\rho\wedge \ext x^\nu$ 反对称，故相消。

\begin{example}[电磁学就是外微分]
$A=A_\mu \ext x^\mu$ 是 1-形式，
\begin{equation}
  F=\ext A=\partial_\mu A_\nu\,\ext x^\mu\wedge \ext x^\nu
  =\frac12\big(\partial_\mu A_\nu-\partial_\nu A_\mu\big)
  \ext x^\mu\wedge \ext x^\nu
  =\frac12 F_{\mu\nu}\ext x^\mu\wedge \ext x^\nu .
\end{equation}
于是 Bianchi 恒等式\textbf{自动成立}：
\begin{equation}
  \ext F=\ext^2A=0
  \qquad\Longleftrightarrow\qquad
  \partial_{[\rho}F_{\mu\nu]}=0 .
  \label{eq:bianchi}
\end{equation}
规范变换 $A\to A+\ext\lambda$ 使 $F$ 不变，同样因 $\ext^2=0$。
\end{example}

\begin{definition}[闭形式与恰当形式]
$\ext\omega=0$ 称 $\omega$ \textbf{闭}（\en{closed}）；
$\omega=\ext\eta$ 称 $\omega$ \textbf{恰当}（\en{exact}）。
由 $\ext^2=0$，恰当 $\Rightarrow$ 闭；反之不一定
（差别由拓扑刻画，见 \eqref{eq:deRham}）。
\end{definition}

\subsubsection{Stokes 定理}

\begin{theorem}[Stokes]
设 $\Mfd_{p+1}$ 是 $(p+1)$ 维带边流形，边界 $\partial\Mfd_{p+1}$，
$\omega$ 为 $p$-形式，则
\begin{equation}
  \boxed{\ \int_{\Mfd_{p+1}}\ext\omega
  =\int_{\partial\Mfd_{p+1}}\omega\ }
  \label{eq:stokes}
\end{equation}
\end{theorem}

\begin{keypoint}
\label{kp:stokes}
\textbf{Stokes 定理是本课程使用频率最高的工具。}
广义对称性中几乎每一次“算符可以自由形变”的论证，
都是 \eqref{eq:stokes} 加上“被积形式是闭的”这两件事的组合。
请把这句话记牢：
\begin{center}
\emph{闭形式在同调等价的子流形上积分值相同。}
\end{center}
\end{keypoint}

\begin{proof}[要点（同调不变性）]
设 $\Ssurf$ 与 $\Ssurf'$ 是两个 $p$ 维闭子流形，
且它们“同调”（\en{homologous}），即存在 $(p+1)$ 维 $V$ 使
$\partial V=\Ssurf-\Ssurf'$。若 $\ext\omega=0$，则
\begin{equation}
  \int_\Ssurf\omega-\int_{\Ssurf'}\omega
  =\int_{\partial V}\omega
  =\int_V \ext\omega=0 .
  \label{eq:homologyinv}
\end{equation}
\end{proof}

\subsubsection{Hodge 星号与守恒律}

\begin{definition}[Hodge 星号]
在 $d$ 维带度规流形上，$\hs:\Omega^p\to\Omega^{d-p}$，
\begin{equation}
  (\hs\omega)_{\mu_1\cdots\mu_{d-p}}
  =\frac{1}{p!}\,\sqrt{|g|}\;
  \epsilon_{\mu_1\cdots\mu_{d-p}\nu_1\cdots\nu_p}\,
  \omega^{\nu_1\cdots\nu_p} .
  \label{eq:hodge}
\end{equation}
在欧氏号差下 $\hs\hs=(-1)^{p(d-p)}$；洛伦兹号差下多一个整体负号。
\end{definition}

现在是关键的翻译。设 $j=j_\mu \ext x^\mu$ 为 1-形式电流，
则 $\hs j$ 是 $(d-1)$-形式，$\ext\hs j$ 是 $d$-形式，正比于体积元：

\begin{proposition}[守恒律的形式语言]
\begin{equation}
  \boxed{\ \partial_\mu j^\mu=0
  \qquad\Longleftrightarrow\qquad
  \ext\hs j=0\ }
  \label{eq:conservation}
\end{equation}
\end{proposition}

\begin{proof}
由 \eqref{eq:hodge}，$(\hs j)_{\mu_1\cdots\mu_{d-1}}
=\sqrt{|g|}\,\epsilon_{\mu_1\cdots\mu_{d-1}\nu}j^\nu$。
取外微分并注意 $d$-形式空间是一维的（正比于
$\ext^dx$），得
\begin{equation}
  \ext\hs j=\big(\partial_\mu j^\mu\big)\sqrt{|g|}\;
  \ext x^1\wedge\cdots\wedge \ext x^d
  =\big(\partial_\mu j^\mu\big)\,\Vol_d .
\end{equation}
故 $\ext\hs j=0$ 当且仅当 $\partial_\mu j^\mu=0$。
\end{proof}

同理，对反对称的 $(q+1)$ 阶张量电流
$J^{\mu_1\cdots\mu_{q+1}}$（即 $(q+1)$-形式 $J$）：
\begin{equation}
  \partial_{\mu_1}J^{\mu_1\cdots\mu_{q+1}}=0
  \qquad\Longleftrightarrow\qquad
  \ext\hs J=0 ,
  \qquad
  \hs J\in\Omega^{d-q-1} .
  \label{eq:conservation-q}
\end{equation}
\textbf{请记住 $\hs J$ 的次数是 $d-q-1$}：
它决定了对称算符支撑在多少维的子流形上（第 \ref{sec:pform} 节）。

\subsubsection{de Rham 上同调与连接数}

\begin{definition}[de Rham 上同调]
\begin{equation}
  H^p(\Mfd,\RR)
  =\frac{\{\text{闭 }p\text{-形式}\}}{\{\text{恰当 }p\text{-形式}\}} .
  \label{eq:deRham}
\end{equation}
它测量“闭但不恰当”的程度，是纯拓扑量。
\end{definition}

\begin{example}[$S^2$ 上的磁通]
球面 $S^2$ 上 $H^2(S^2)\cong\RR\ne0$。
磁单极子场强 $F$ 满足 $\ext F=0$ 但\textbf{不能}整体写成 $\ext A$，
其“通量” $\frac{1}{2\pi}\oint_{S^2}F=m\in\ZZ$ 就是拓扑数
（Dirac 量子化）。这个例子将在第 \ref{subsec:tHooft-line} 小节直接使用。
\end{example}

\begin{definition}[连接数]
\label{def:link}
设 $\Mfd_p$ 与 $\Mfd_q$ 是 $d$ 维时空中两个\textbf{不相交}的闭子流形，
且
\begin{equation}
  \boxed{\ p+q+1=d\ }
  \label{eq:linkdim}
\end{equation}
则可定义整数\textbf{连接数}（\en{linking number}）$\Link(\Mfd_p,\Mfd_q)$：
取 $(p+1)$ 维 $V$ 使 $\partial V=\Mfd_p$，
则 $\Link=$ $V$ 与 $\Mfd_q$ 的（带符号）相交数。
\end{definition}

\begin{example}[$d=3$：两条圈]
$p=q=1$，$1+1+1=3$ $\checkmark$。这就是通常绳结理论里两条闭曲线的连接数。
\end{example}

\begin{example}[$d=4$：2 维面与 1 维圈]
$p=2,q=1$，$2+1+1=4$ $\checkmark$。
一个 2 维球面 $\Ssurf_2$ 可以“套住”一条曲线 $C$。
这正是第 \ref{subsec:maxwell} 小节中对称算符与 Wilson 线的关系。
\end{example}

\begin{keypoint}
\label{kp:linkdim}
条件 \eqref{eq:linkdim} 是全课程的\textbf{维数记账法则}。
遇到任何“对称算符作用在带荷对象上”的陈述，
先用它检查维数是否配得上：配不上就一定记错了。
\end{keypoint}

\paragraph{$\delta$-函数形式（Poincaré 对偶）。}
对 $d$ 维时空中的 $k$ 维闭子流形 $\Mfd_k$，
定义 $(d-k)$-形式 $\delta_{\Mfd_k}$，使得对任意 $k$-形式 $\omega$
\begin{equation}
  \int_{\Mfd_d}\delta_{\Mfd_k}\wedge\omega
  =\int_{\Mfd_k}\omega .
  \label{eq:poincare}
\end{equation}
它是“支撑在 $\Mfd_k$ 上的 $\delta$ 函数”。
这个记号能把“线算符插入”写进运动方程，见 \eqref{eq:modifiedEOM}。

\begin{exercise}
\label{ex:forms}
(a) 在 $d=4$ 中验证 $F\wedge F$ 是 4-形式，
并证明 $F\wedge F=\ext(A\wedge F)$，
即 $\int F\wedge F$ 是拓扑项（$\theta$ 项）。\\
(b) 用 \eqref{eq:hodge} 验证 $d=4$ 中
$\hs(\ext x^0\wedge \ext x^1)\propto \ext x^2\wedge \ext x^3$。\\
(c) 对 $q$-form 电流，用 \eqref{eq:linkdim} 验证：
若对称算符支撑在 $(d-q-1)$ 维流形上，
则带荷对象必为 $q$ 维。
\end{exercise}

% =====================================================================
\section{普通对称性：连续／离散、全局／规范}
\label{sec:ordinary}

在推广之前，必须把“普通对称性”的三组区分讲得非常干净，
因为广义对称性的每一步推广都精确地对应着放松其中某一条。

\subsection{对称性的定义}

\begin{definition}[经典对称性]
一个场变换 $\phi^a\to\phi'^a=\phi^a+\delta\phi^a$ 称为作用量 $S$ 的
\textbf{对称性}，若
\begin{equation}
  S[\phi']=S[\phi]
  \qquad\text{或更一般地}\qquad
  \delta\Lag=\partial_\mu K^\mu
  \label{eq:symdef}
\end{equation}
即拉氏量最多变化一个全导数（此时作用量在合适边界条件下不变）。
\end{definition}

\begin{warning}
量子层面上还要求\textbf{路径积分测度}不变；
测度不变性的破坏就是\textbf{反常}（\en{anomaly}）。
第 \ref{sec:ssb} 节会区分两类反常：
“阻碍规范化”的 't Hooft 反常，与“真正破坏守恒律”的 ABJ 型反常。
\end{warning}

\subsection{连续 vs 离散}

\begin{itemize}
  \item \textbf{连续对称性}（\en{continuous symmetry}）：
        变换由连续参数标记，如 $\phi\to e^{i\alpha}\phi$。
        存在无穷小形式 $\delta\phi=i\alpha\phi$，
        因此\textbf{有 Noether 流}（第 \ref{sec:noether} 节）。
  \item \textbf{离散对称性}（\en{discrete symmetry}）：
        如 $\varphi\to-\varphi$（$\ZZ_2$）、时间反演 $T$、
        空间反射 $P$、电荷共轭 $C$。\textbf{没有}连续参数，
        因此\textbf{没有守恒流}，只有守恒的算符 $U$（$U^2=1$ 等）。
\end{itemize}

\begin{keypoint}
\label{kp:discrete}
传统教科书对离散对称性的处理是“瘸腿”的：
有对称算符 $U$ 却没有流，于是所有以 $j^\mu$ 为中心的技术
（Ward 恒等式、荷代数）都用不上。

\textbf{广义对称性的第一个好处就是修好了这条腿}：
把对称性定义为\textbf{拓扑算符}而不是流，
连续与离散就被统一处理了——连续情形的算符恰好可以写成
$\exp(i\alpha\int\hs j)$，离散情形则不能，但两者都是拓扑算符。
这也是为什么本课程能同时讨论 $\ZZ_2$ 与 $\Uone$。
\end{keypoint}

\subsection{全局 vs 规范：本课程最重要的区分}
\label{subsec:global-vs-gauge}

\begin{definition}[全局与规范]
若变换参数 $\alpha$ 是\textbf{常数}（不依赖时空点），
称\textbf{全局对称性}（\en{global symmetry}）；
若参数可取任意函数 $\alpha(x)$ 而作用量仍不变
（通常需引入规范场），称\textbf{规范对称性}（\en{gauge symmetry}）。
\end{definition}

\begin{warning}[必须彻底接受的观念]
\label{warn:gauge}
\textbf{规范对称性不是对称性，而是描述的冗余}
（\en{redundancy of description}）。理由：
\begin{itemize}
  \item 规范变换\textbf{不改变物理态}：$|\psi\rangle$ 与其规范变换像
        是 Hilbert 空间中\emph{同一个}矢量；物理 Hilbert 空间由
        规范不变态构成（Gauss 律约束）。
  \item 因此规范“对称性”\textbf{不能自发破缺}，
        也没有相应的选择定则。所谓“Higgs 机制破坏规范对称性”
        是一种历史遗留的说法，严格讲被破坏的是全局部分。
  \item \textbf{全局}对称性才有物理后果：
        选择定则、Goldstone 玻色子、't Hooft 反常约束。
\end{itemize}
本课程标题里的“广义对称性”全称是\textbf{广义全局对称性}
（\en{generalized global symmetry}）——中间那个 “global” 不是装饰。
\end{warning}

\begin{keypoint}
\label{kp:gauge-produces}
虽然规范对称性本身不是对称性，
它却是\textbf{产生}高阶形式对称性的主要机制：
\begin{center}
\emph{规范场的存在使得线／面算符（Wilson 线等）成为自然的观测量，
而作用在这些拓展算符上的对称性就是高阶形式对称性。}
\end{center}
这就是为什么第 \ref{sec:examples} 节的核心例子全是规范理论。
\end{keypoint}

\begin{example}[对比表]
考虑复标量 $\phi$，电荷 1。
\begin{center}
\begin{tabular}{@{}lll@{}}
\toprule
 & 全局 $\Uone$ & 规范 $\Uone$ \\
\midrule
变换 & $\phi\to e^{i\alpha}\phi$，$\alpha$ 常数
     & $\phi\to e^{i\alpha(x)}\phi$，$A_\mu\to A_\mu+\partial_\mu\alpha$ \\
是否需引入 $A_\mu$ & 否 & 是 \\
守恒流 & $j^\mu$，物理可测 & $j^\mu$ 成为 Gauss 律约束 \\
可否自发破缺 & 可以（Goldstone） & 不可（只是冗余）\\
物理后果 & 选择定则、粒子数守恒 & 定义相互作用、无独立后果 \\
\bottomrule
\end{tabular}
\end{center}
\end{example}

\begin{exercise}
说明为什么“规范对称性不能自发破缺”与“超导体中 $\Uone$ 被破坏”
这两句话并不矛盾。（提示：超导中被破坏的是电子数的\emph{全局}
$\Uone$；而规范场获得质量的现象在广义对称性语言中可描述为
1-form 对称性的行为，见第 \ref{subsec:confinement} 小节。）
\end{exercise}

% =====================================================================
\section{Noether 定理完整推导}
\label{sec:noether}

本节给出完整推导：变分 $\to$ 守恒流 $\to$ 守恒荷 $\to$
（关键一步）\textbf{拓扑算符}。前三步是标准内容，
第四步是通往广义对称性的门。

\subsection{定理与推导}

\begin{theorem}[Noether]
\label{thm:noether}
设 $S=\int \ext^dx\,\Lag(\phi^a,\partial_\mu\phi^a)$ 在无穷小变换
\begin{equation}
  \phi^a\longrightarrow\phi^a+\epsilon\,\Delta^a(\phi,\partial\phi)
  \label{eq:infsym}
\end{equation}
下满足 $\delta\Lag=\epsilon\,\partial_\mu K^\mu$（$\epsilon$ 为常数）。
则在运动方程成立时（\en{on-shell}），电流
\begin{equation}
  \boxed{\ j^\mu=\PD{\Lag}{(\partial_\mu\phi^a)}\,\Delta^a-K^\mu\ }
  \label{eq:noethercurrent}
\end{equation}
满足守恒律 $\partial_\mu j^\mu=0$。
\end{theorem}

\begin{proof}
\textbf{第 1 步（变分）。}直接对 $\Lag$ 做变分，
注意 $\delta(\partial_\mu\phi^a)=\partial_\mu(\delta\phi^a)$：
\begin{equation}
  \delta\Lag
  =\PD{\Lag}{\phi^a}\,\delta\phi^a
  +\PD{\Lag}{(\partial_\mu\phi^a)}\,\partial_\mu\big(\delta\phi^a\big),
  \qquad \delta\phi^a=\epsilon\Delta^a .
  \label{eq:noether1}
\end{equation}

\textbf{第 2 步（用运动方程消去 $\partial\Lag/\partial\phi$）。}
由 Euler--Lagrange 方程 \eqref{eq:EL}，
\begin{equation}
  \PD{\Lag}{\phi^a}
  =\partial_\mu\!\left(\PD{\Lag}{(\partial_\mu\phi^a)}\right).
  \label{eq:noether2}
\end{equation}
代入 \eqref{eq:noether1}：
\begin{align}
  \delta\Lag
  &=\partial_\mu\!\left(\PD{\Lag}{(\partial_\mu\phi^a)}\right)\epsilon\Delta^a
   +\PD{\Lag}{(\partial_\mu\phi^a)}\,\partial_\mu\big(\epsilon\Delta^a\big)
  \notag\\[2pt]
  &=\epsilon\,\partial_\mu\!\left[
    \PD{\Lag}{(\partial_\mu\phi^a)}\,\Delta^a\right],
  \label{eq:noether3}
\end{align}
最后一步就是乘积法则的逆用（这里用到 $\epsilon$ 是常数）。

\textbf{第 3 步（与假设比较）。}
按对称性假设 $\delta\Lag=\epsilon\,\partial_\mu K^\mu$，
与 \eqref{eq:noether3} 相减：
\begin{equation}
  \epsilon\,\partial_\mu\!\left[
  \PD{\Lag}{(\partial_\mu\phi^a)}\Delta^a-K^\mu\right]=0 .
\end{equation}
因 $\epsilon\ne0$ 任意，方括号的散度为零，即 \eqref{eq:noethercurrent}。
\end{proof}

\begin{remark}[为什么必须“在壳”]
第 2 步用了运动方程。这是本质的：
\textbf{Noether 守恒律是运动方程的推论}，
而不是恒等式。对比 \eqref{eq:bianchi} 的 Bianchi 恒等式
$\ext F=0$——它\emph{不}依赖运动方程，是恒等式。
第 \ref{subsec:maxwell} 小节会看到，这两类来源
分别给出 Maxwell 理论的\textbf{电}与\textbf{磁} 1-form 对称性，
而“一个来自运动方程、一个来自恒等式”这个区别正是
它们在加入物质场后命运不同的原因。
\end{remark}

\subsection{从流到荷}
\label{subsec:charge}

\begin{proposition}[守恒荷]
定义在固定时刻 $t$ 的空间切片上的积分
\begin{equation}
  Q(t)=\int_{\RR^{d-1}}\ext^{d-1}x\;j^0(t,\bm x),
  \label{eq:chargedef}
\end{equation}
若场在空间无穷远处足够快衰减，则 $\dfrac{\ext Q}{\ext t}=0$。
\end{proposition}

\begin{proof}
\begin{equation}
  \frac{\ext Q}{\ext t}
  =\int \ext^{d-1}x\;\partial_0 j^0
  \overset{\partial_\mu j^\mu=0}{=}
  -\int \ext^{d-1}x\;\partial_i j^i
  \overset{\text{Stokes}}{=}
  -\oint_{S^{d-2}_\infty}\ext S_i\;j^i=0 .
\end{equation}
\end{proof}

\begin{warning}
“衰减足够快”不是小事。若允许场在无穷远不衰减，
$Q$ 可以不守恒，或者出现\textbf{额外}的守恒荷。
本课程后面会看到，正是对“边界／无穷远行为”的仔细处理
区分了不同的全局结构。
\end{warning}

\subsection{两个完整算例}
\label{subsec:two-examples}

\begin{example}[复标量的 $\Uone$ 流]
\label{ex:u1current}
取例 \ref{ex:threelag}(2)：
$\Lag=-\partial_\mu\phi^*\partial^\mu\phi-V(|\phi|^2)$。
对称变换 $\phi\to e^{i\alpha}\phi$，无穷小
\begin{equation}
  \delta\phi=i\alpha\phi,\qquad \delta\phi^*=-i\alpha\phi^* .
\end{equation}
势项只依赖 $|\phi|^2$ 故不变；动能项也不变，所以 $K^\mu=0$。
计算导数：
\begin{equation}
  \PD{\Lag}{(\partial_\mu\phi)}=-\partial^\mu\phi^*,
  \qquad
  \PD{\Lag}{(\partial_\mu\phi^*)}=-\partial^\mu\phi .
\end{equation}
代入 \eqref{eq:noethercurrent}（把 $\alpha$ 提出，$\Delta=i\phi$，
$\Delta^*=-i\phi^*$）：
\begin{align}
  j^\mu
  &=\big(-\partial^\mu\phi^*\big)\big(i\phi\big)
   +\big(-\partial^\mu\phi\big)\big(-i\phi^*\big)
  \notag\\
  &=i\left(\phi^*\partial^\mu\phi-\phi\,\partial^\mu\phi^*\right).
  \label{eq:u1current}
\end{align}
验证守恒：
\begin{equation}
  \partial_\mu j^\mu
  =i\left(\phi^*\Box\phi-\phi\,\Box\phi^*\right)
  \overset{\text{EOM}}{=}
  i\left(\phi^*\,\phi\,V'-\phi\,\phi^*\,V'\right)=0 . \checkmark
\end{equation}
（用了运动方程 $\Box\phi=V'(|\phi|^2)\phi$ 及其复共轭。）
\end{example}

\begin{example}[平移与能量--动量张量]
\label{ex:stress}
取变换为时空平移 $x^\mu\to x^\mu-\epsilon^\nu\delta^\mu_\nu$，
即 $\delta\phi^a=\epsilon^\nu\partial_\nu\phi^a$。
此时 $\Lag$ 本身也平移：$\delta\Lag=\epsilon^\nu\partial_\nu\Lag
=\epsilon^\nu\partial_\mu(\delta^\mu_\nu\Lag)$，
故 $K^\mu{}_\nu=\delta^\mu_\nu\Lag$。代入 \eqref{eq:noethercurrent}：
\begin{equation}
  \boxed{\ T^\mu{}_\nu
  =\PD{\Lag}{(\partial_\mu\phi^a)}\,\partial_\nu\phi^a
  -\delta^\mu_\nu\,\Lag\ },
  \qquad
  \partial_\mu T^\mu{}_\nu=0 .
  \label{eq:stresstensor}
\end{equation}
每个 $\nu$ 给一个守恒流；对应的荷是能量与动量
$P_\nu=\int \ext^{d-1}x\,T^0{}_\nu$。
\end{example}

\begin{exercise}
对实标量 $\Lag=-\tfrac12(\partial\varphi)^2-V(\varphi)$，
证明 $\ZZ_2$ 对称性 $\varphi\to-\varphi$ \textbf{不}给出 Noether 流
（提示：$\Delta$ 必须是无穷小的，而 $\varphi\to-\varphi$ 无法连续地
连到恒等变换）。这再次说明要点 \ref{kp:discrete} 的必要性。
\end{exercise}

\subsection{荷生成对称变换（正则验证）}
\label{subsec:generates}

守恒荷不只是个数，它是\textbf{生成元}。这里用例
\ref{ex:u1current} 做完整验证。

自由复标量（$V=m^2|\phi|^2$）的共轭动量
\begin{equation}
  \pi=\PD{\Lag}{\dot\phi}=\dot\phi^*,
  \qquad
  \pi^*=\PD{\Lag}{\dot\phi^*}=\dot\phi .
\end{equation}
由 \eqref{eq:u1current} 取 $\mu=0$（号差 $(-,+,+,+)$ 下
$\partial^0=-\partial_0$，此处把整体符号吸收进 $Q$ 的定义，
最终结论不依赖该约定）：
\begin{equation}
  j^0=i\big(\phi\,\pi-\phi^*\pi^*\big),
  \qquad
  Q=\int \ext^{d-1}x\;i\big(\phi\,\pi-\phi^*\pi^*\big).
  \label{eq:Qcanonical}
\end{equation}
用正则关系 \eqref{eq:canonical}，
$[\phi(\bm x),\pi(\bm y)]=i\delta^{d-1}(\bm x-\bm y)$：
\begin{align}
  \big[Q,\phi(\bm x)\big]
  &=\int \ext^{d-1}y\;i\,\phi(\bm y)\,
    \big[\pi(\bm y),\phi(\bm x)\big]
  \notag\\
  &=\int \ext^{d-1}y\;i\,\phi(\bm y)\,
    \big(-i\,\delta^{d-1}(\bm x-\bm y)\big)
  \;=\;\phi(\bm x).
  \label{eq:Qphi}
\end{align}
（$\phi^*\pi^*$ 项与 $\phi$ 交换为零。）于是对有限参数
\begin{equation}
  \boxed{\ \Uop_\alpha=e^{i\alpha Q}:
  \qquad
  \Uop_\alpha\,\phi(x)\,\Uop_\alpha^{-1}
  =\phi+i\alpha[Q,\phi]+\cdots
  =e^{i\alpha}\phi(x)\ }
  \label{eq:Uacts}
\end{equation}
恰好是我们出发的对称变换。$\checkmark$

\begin{keypoint}
\label{kp:U-operator}
\eqref{eq:Uacts} 定义的 $\Uop_\alpha=e^{i\alpha Q}$ 就是\textbf{对称算符}
（\en{symmetry operator}）。注意它满足
\begin{equation}
  \Uop_\alpha\,\Uop_\beta=\Uop_{\alpha+\beta},
  \qquad
  \Uop_\alpha^{-1}=\Uop_{-\alpha},
  \label{eq:groupfusion}
\end{equation}
即\textbf{可逆}，且乘法就是群 $\Uone$ 的乘法。
广义对称性放松的正是 \eqref{eq:groupfusion}：
若只要求算符相乘时闭合成某个代数而不要求有逆，
就得到不可逆对称性。
\end{keypoint}

\subsection{Ward 恒等式：荷如何“测量”算符}
\label{subsec:ward}

在关联函数层面，对称性表现为 Ward 恒等式。
设 $\Op(y)$ 是带荷 $q$ 的局域算符，即
\begin{equation}
  \Uop_\alpha\,\Op(y)\,\Uop_\alpha^{-1}=e^{iq\alpha}\,\Op(y),
  \qquad\text{无穷小地}\qquad
  \delta\Op=i q\alpha\,\Op .
\end{equation}

\paragraph{推导。}在路径积分 \eqref{eq:pathintegral} 中做一个\textbf{局域}的
变换 $\phi\to\phi+\epsilon(x)\Delta$，其中参数 $\epsilon$ 现在依赖 $x$。
因为常数 $\epsilon$ 时 $\delta S=0$，一般 $\epsilon(x)$ 时 $\delta S$ 必须
正比于 $\partial_\mu\epsilon$：
\begin{equation}
  \delta S=-\int \ext^dx\;\epsilon(x)\,\partial_\mu j^\mu ,
  \label{eq:deltaSlocal}
\end{equation}
（把 $\int\partial_\mu\epsilon\,j^\mu$ 分部积分而得）。
路径积分对场的重命名不变，故
\begin{equation}
  0=\delta\Big\langle \Op(y)\Big\rangle
  =\Big\langle i\,\delta S\;\Op(y)\Big\rangle
   +\Big\langle \delta\Op(y)\Big\rangle .
\end{equation}
代入 \eqref{eq:deltaSlocal} 与 $\delta\Op(y)=i q\,\epsilon(y)\Op(y)$，
并对任意 $\epsilon(x)$ 成立，得

\begin{proposition}[Ward 恒等式]
\begin{equation}
  \boxed{\
  \partial_\mu\big\langle j^\mu(x)\,\Op(y)\big\rangle
  = q\;\delta^{d}(x-y)\,\big\langle \Op(y)\big\rangle \ }
  \label{eq:ward}
\end{equation}
（整体符号依赖 $j$ 与 $Q$ 的符号约定；结构是关键。）
\end{proposition}

\paragraph{最重要的推论。}把 \eqref{eq:ward} 在一个包住 $y$ 的
$d$ 维区域 $V$ 上积分，并用 Stokes 定理：
\begin{equation}
  \oint_{\partial V}\big\langle \hs j\;\Op(y)\big\rangle
  =\int_V \big\langle \ext\hs j\;\Op(y)\big\rangle
  = q\,\big\langle\Op(y)\big\rangle
  \qquad (y\in V).
  \label{eq:chargebylinking}
\end{equation}

\begin{keypoint}
\label{kp:measure}
\eqref{eq:chargebylinking} 是本讲义\textbf{第一个真正重要的结论}：
\begin{center}
\emph{把对称算符支撑在一个包围算符 $\Op$ 的闭曲面上，
测得的就是 $\Op$ 的荷。}
\end{center}
换句话说，\textbf{“带荷”是一个关于“包围”（拓扑连接）的陈述}，
而不是关于“在同一点”的陈述。
一旦这样理解，把“包围”换成“连接”就自然地推广到
高维带荷对象——这就是 $p$-form 对称性。
\end{keypoint}

\subsection{关键改写：\texorpdfstring{$Q$}{Q} 是拓扑算符}
\label{subsec:topological-rewrite}

现在做整门课的枢纽一步。把 \eqref{eq:chargedef} 的荷用微分形式重写：

\begin{definition}[曲面上的荷]
对 $d$ 维时空中任意闭的 $(d-1)$ 维子流形 $\Ssurf_{d-1}$，定义
\begin{equation}
  \boxed{\ Q(\Ssurf_{d-1})=\int_{\Ssurf_{d-1}}\hs j ,
  \qquad
  \Uop_\alpha(\Ssurf_{d-1})=\exp\!\Big[i\alpha\!\int_{\Ssurf_{d-1}}\hs j\Big]\ }
  \label{eq:Qsurface}
\end{equation}
\end{definition}

取 $\Ssurf_{d-1}=\{t=\text{const}\}$ 的空间切片就回到
\eqref{eq:chargedef}（因为 $\hs j$ 在等时面上的分量正是 $j^0$）。
但 \eqref{eq:Qsurface} 允许 $\Ssurf$ 是\textbf{任意}闭曲面——
甚至是弯曲的、不等时的、多个连通分支的。

\begin{theorem}[对称算符是拓扑的]
\label{thm:topological}
设 $\ext\hs j=0$（即 $\partial_\mu j^\mu=0$）。
若 $\Ssurf$ 与 $\Ssurf'$ 同调，即存在 $d$ 维区域 $V$ 使
$\partial V=\Ssurf-\Ssurf'$，且 $V$ 内没有带荷算符插入，则
\begin{equation}
  Q(\Ssurf)=Q(\Ssurf'),
  \qquad
  \Uop_\alpha(\Ssurf)=\Uop_\alpha(\Ssurf') .
  \label{eq:topological}
\end{equation}
\end{theorem}

\begin{proof}
直接用 Stokes 定理 \eqref{eq:stokes} 与守恒律：
\begin{equation}
  Q(\Ssurf)-Q(\Ssurf')
  =\int_{\partial V}\hs j
  =\int_V \ext\hs j
  =0 .
\end{equation}
\end{proof}

\begin{keypoint}
\label{kp:pivot}
定理 \ref{thm:topological} 就是标语 \eqref{eq:slogan}。请体会它的分量：
\begin{enumerate}[label=(\arabic*)]
  \item \textbf{“守恒”与“拓扑”是同一件事的两种说法。}
        $\partial_\mu j^\mu=0$ $\Leftrightarrow$
        $\Uop_\alpha(\Ssurf)$ 在连续形变下不变。
  \item \textbf{对称算符不必放在等时面上。}
        它是一个支撑在\emph{某个子流形}上的算符，
        这件事在传统表述里完全被“$Q=\int \ext^{d-1}x\,j^0$”掩盖了。
  \item \textbf{“不变”有前提}：形变过程中不能扫过带荷算符。
        若扫过了，$Q$ 就跳变一个荷量——这正是
        \eqref{eq:chargebylinking} 的内容，也是对称性作用的机制。
  \item \textbf{定义里已经不再需要 $j$ 了。}
        我们真正用到的只是“存在一族支撑在 $(d-1)$ 维闭曲面上的、
        拓扑的、可逆的算符”。把 $j$ 扔掉，
        离散对称性也被涵盖；把维数改掉，$p$-form 对称性就出现了。
\end{enumerate}
\end{keypoint}

\begin{remark}[两种“不变”的对比]
初学时容易混淆两句话：
\begin{itemize}
  \item “$Q$ 不随时间变化”——传统说法，
        对应 $\Ssurf$ 与 $\Ssurf'$ 是两个不同时刻的等时面。
  \item “$Q(\Ssurf)$ 只依赖 $\Ssurf$ 的同调类”——广义说法，
        上一条是它的特例（两个等时面确实同调，
        中间夹的 $V$ 就是这段时空）。
\end{itemize}
所以广义表述\textbf{严格更强}，并没有丢掉任何旧内容。
\end{remark}

\begin{exercise}
\label{ex:pivot}
(a) 在 $d=4$ 中，取 $\Ssurf$ 为半径 $R$ 的球面 $\times$ 时间区间的边界，
说明 \eqref{eq:topological} 如何同时包含“电荷守恒”与“电荷与 $R$ 无关”。\\
(b) 若 $V$ 内含有一个荷为 $q$ 的算符 $\Op$，
把定理 \ref{thm:topological} 的证明与 \eqref{eq:chargebylinking} 结合，
写出 $Q(\Ssurf)-Q(\Ssurf')$ 等于多少。\\
(c) 用 \eqref{eq:linkdim} 检查：$(d-1)$ 维对称算符与
$0$ 维局域算符的连接数是否有定义？
（答：$(d-1)+0+1=d$ $\checkmark$。）
\end{exercise}

% =====================================================================
\section{QFT 中的算符与关联函数：局域 vs 拓展}
\label{sec:operators}

要谈“作用在线上的对称性”，先要承认线本身是合法的观测量。
本节把算符按\textbf{支撑集维数}分类。

\subsection{关联函数是理论的定义}

一个 QFT 的全部信息可以编码在所有算符的关联函数中：
\begin{equation}
  \big\langle \Op_1(x_1)\,\Op_2(x_2)\cdots\big\rangle .
\end{equation}
两个理论若所有关联函数相同，就是同一个理论。
所以\textbf{“有哪些算符”是理论定义的一部分}——
这句话在第 \ref{subsec:duality} 小节会变得非常实在：
$\Uone$ 与 $\Uone/\ZZ_N$ 规范理论的\emph{局域}动力学完全相同，
差别只在\textbf{允许哪些线算符存在}。

\subsection{局域算符}

\begin{definition}[局域算符]
\textbf{局域算符}（\en{local operator}）支撑在时空中的一个\textbf{点}上：
$\Op(x)$。它是 $0$ 维算符。
\end{definition}

例子：$\varphi(x)$、$\phi^*\phi(x)$、$T^{\mu\nu}(x)$、
$F_{\mu\nu}F^{\mu\nu}(x)$，以及 2 维 CFT 中的顶点算符
$e^{i n\varphi(x)}$。

\begin{warning}
在规范理论中，$\phi(x)$ 本身\textbf{不是}好的算符（不规范不变）。
规范不变的局域算符必须是如 $|\phi|^2$、$F_{\mu\nu}F^{\mu\nu}$ 之类。
这个限制\textbf{正是}拓展算符登场的原因：
带电荷的信息无法装进局域算符，只能装进线算符。
\end{warning}

\subsection{拓展算符}
\label{subsec:extended}

\begin{definition}[拓展算符]
\textbf{拓展算符}（\en{extended operator}）支撑在时空中一个
$p$ 维子流形上（$p\ge1$），记作 $\Op(\Mfd_p)$。
$p=1$ 称\textbf{线算符}（\en{line operator}），
$p=2$ 称\textbf{面算符}（\en{surface operator}），依此类推。
\end{definition}

\begin{example}[Wilson 线：最重要的线算符]
\label{ex:wilson}
对 $\Uone$ 规范场，沿闭曲线 $C$ 定义
\begin{equation}
  \boxed{\ W_n(C)=\exp\!\left[i\,n\oint_C A\right]
  =\exp\!\left[i\,n\oint_C A_\mu\,\ext x^\mu\right],
  \qquad n\in\ZZ\ }
  \label{eq:wilson}
\end{equation}
\textbf{规范不变性}：在 $A\to A+\ext\lambda$ 下
\begin{equation}
  \oint_C A\;\longrightarrow\;\oint_C A+\oint_C \ext\lambda
  =\oint_C A+\int_{\partial C}\lambda
  =\oint_C A ,
\end{equation}
因为 $C$ 闭合故 $\partial C=\emptyset$。$\checkmark$

\textbf{整数量子化}：若 $\lambda$ 是允许绕圈的（大规范变换），
$\oint_C \ext\lambda\in2\pi\ZZ$，要求 $W$ 单值就需 $n\in\ZZ$。

\textbf{物理意义}：$W_n(C)$ 是一个\textbf{电荷为 $n$ 的无穷重（探针）粒子
沿 $C$ 运动的世界线}。取 $C$ 为时间方向的直线即一个静止的静态电荷。

对非阿贝尔群，需路径有序化并取迹：
\begin{equation}
  W_R(C)=\Tr_R\,\Pexp\!\left[i\oint_C A\right].
  \label{eq:wilsonNA}
\end{equation}
\end{example}

\begin{example}['t Hooft 线]
\label{ex:thooft}
\textbf{'t Hooft 线}（\en{'t Hooft line}）$T_m(C)$ 是\textbf{磁}荷 $m$ 的
探针粒子（磁单极子）的世界线。它不像 \eqref{eq:wilson} 那样有简单的
“$\exp\oint$”表达式，而是通过\textbf{边界条件}定义：
要求在环绕 $C$ 的小球面 $S^2$ 上
\begin{equation}
  \frac{1}{2\pi}\oint_{S^2}F=m\in\ZZ .
  \label{eq:thooftdef}
\end{equation}
即“挖掉 $C$ 并规定其周围有 $m$ 个单位磁通”。
在对偶描述中（$F\leftrightarrow\hs F$）它变成 Wilson 线。
\end{example}

\begin{example}[面算符与其他]
2-维面算符的例子：
\begin{itemize}
  \item $\exp\big[i\alpha\int_{\Ssurf_2}\hs F/e^2\big]$——
        第 \ref{subsec:maxwell} 小节的\textbf{对称算符}本身。
  \item 2 维中的\textbf{拓扑缺陷线}（\en{topological defect line}, TDL）——
        第 \ref{subsec:ising} 小节的 Ising 例子。
  \item 弦的世界面、畴壁（\en{domain wall}）。
\end{itemize}
\end{example}

\subsection{“真算符” vs “可端接的算符”}
\label{subsec:genuine}

一个重要而微妙的区分：

\begin{definition}[真算符与非真算符]
一个线算符若能\textbf{独立}定义在任意闭曲线上，
称\textbf{真线算符}（\en{genuine line operator}）。
若它必须作为某个面算符的边界出现
（即“线只能是面的边”），则称\textbf{非真的}
（\en{non-genuine}）。
\end{definition}

\begin{example}[被屏蔽的 Wilson 线]
\label{ex:screening}
若理论中存在\textbf{动力学的}电荷 1 的物质场 $\psi$，
则 $W_1$ 可以“断开”：一段开放的 Wilson 线
$\bar\psi(x)\,e^{i\int_x^y A}\,\psi(y)$ 是规范不变的局域算符对。
于是 $W_1(C)$ 不再携带守恒的荷——它可以被物质屏蔽
（\en{screened}）。
\textbf{这就是电 1-form 对称性被动力学电荷破坏的机制}
（第 \ref{subsec:matter} 小节量化）。
\end{example}

\begin{keypoint}
\label{kp:genuine}
判断一个理论有没有 1-form 对称性，操作上就是问：
\begin{center}
\emph{存在不能被任何动力学物质端接或屏蔽的线算符吗？}
\end{center}
有 $\Rightarrow$ 它们按某个群分类，那个群就是 1-form 对称群。
\end{keypoint}

% =====================================================================
\section{拓扑算符的物理图像}
\label{sec:topological}

第 \ref{subsec:topological-rewrite} 小节已经证明对称算符是拓扑的。
本节反过来，把“拓扑算符”当作\textbf{出发点}，
建立完整的物理图像。

\subsection{定义与判据}

\begin{definition}[拓扑算符]
\label{def:topological}
支撑在闭子流形 $\Mfd_k$ 上的算符 $\Uop(\Mfd_k)$ 称为\textbf{拓扑的}
（\en{topological}），若在\textbf{任何}关联函数中，
连续形变 $\Mfd_k\to\Mfd_k'$ 都不改变结果，
只要形变过程中 $\Mfd_k$ \textbf{不扫过其他算符插入点}：
\begin{equation}
  \big\langle \Uop(\Mfd_k)\;\Op_1\cdots\Op_n\big\rangle
  =\big\langle \Uop(\Mfd_k')\;\Op_1\cdots\Op_n\big\rangle .
  \label{eq:topodef}
\end{equation}
\end{definition}

\begin{remark}[等价的三种判据]
以下三条等价，实际计算中常混用：
\begin{enumerate}[label=(\roman*)]
  \item 关联函数在形变下不变（\eqref{eq:topodef}，最本质）。
  \item 该算符与能量--动量张量“解耦”：
        $\Uop$ 不依赖度规，即 $\delta\Uop/\delta g_{\mu\nu}=0$。
        直觉：不依赖度规 $\Rightarrow$ 不知道距离与形状 $\Rightarrow$ 拓扑。
  \item 对连续对称性，被积形式是闭的：$\ext\hs J=0$
        （由定理 \ref{thm:topological}）。
\end{enumerate}
\end{remark}

\begin{example}[对照：不是拓扑的算符]
$\exp\big[i\oint_C A\big]$（Wilson 线）\textbf{不是}拓扑算符：
它的期望值依赖 $C$ 的面积或周长（第 \ref{subsec:confinement} 小节）。
而 $\exp\big[i\alpha\oint_{\Ssurf_2}\hs F/e^2\big]$ \textbf{是}拓扑算符
（因为 $\ext\hs F=0$）。
\textbf{这个对照极其重要}：Wilson 线是\emph{带荷对象}，
面算符是\emph{对称算符}，两者角色完全不同，不要混为一谈。
\end{example}

\subsection{融合：拓扑算符的乘法}
\label{subsec:fusion}

把两个拓扑算符支撑在同一个（或无限接近的）子流形上，
它们的乘积仍是拓扑算符，可以按某种规则重新展开：

\begin{definition}[融合规则]
\begin{equation}
  \Uop_a(\Mfd_k)\times \Uop_b(\Mfd_k)
  =\sum_c N_{ab}^{\ \ c}\;\Uop_c(\Mfd_k),
  \qquad N_{ab}^{\ \ c}\in\ZZ_{\ge0}
  \label{eq:fusion}
\end{equation}
称\textbf{融合规则}（\en{fusion rule}）。
\end{definition}

两种情形：

\begin{itemize}
  \item \textbf{可逆（群）情形}：右边只有一项且系数为 1，
        $\Uop_a\times\Uop_b=\Uop_{a\cdot b}$。
        此时 $\{\Uop_a\}$ 构成群，存在 $\Uop_{a^{-1}}$ 使
        $\Uop_a\times\Uop_{a^{-1}}=\Uop_{\mathbf 1}$。
        这是普通对称性，如 \eqref{eq:groupfusion}。
  \item \textbf{不可逆情形}：右边是多项之和，
        于是\textbf{没有}逆元。此时称
        \textbf{不可逆对称性}（\en{non-invertible symmetry}）。
        典型例子（第 \ref{subsec:ising} 小节）：
        \begin{equation}
          \mathcal N\times\mathcal N=\mathbf 1+\eta .
          \label{eq:isingfusionpreview}
        \end{equation}
\end{itemize}

\begin{keypoint}
\label{kp:whysymmetry}
为什么把不可逆的东西也叫“对称性”？因为对称性的\textbf{实用价值}
——给出关联函数的\textbf{选择定则与约束}——只需要
“存在拓扑算符”这一条，并不需要可逆性。
拓扑性保证算符可以自由移动、穿过、包围其他算符，
从而产生恒等式；有没有逆元并不影响这一点。
\end{keypoint}

\subsection{作用机制：穿过与包围}
\label{subsec:action}

拓扑算符怎样“作用”在别的算符上？有两个等价图像。

\paragraph{图像一：包围（测量荷）。}
把 $\Uop_\alpha(\Ssurf_{d-1})$ 收缩成一个包围 $\Op(y)$ 的小球面。
由 \eqref{eq:chargebylinking}，
\begin{equation}
  \Uop_\alpha\big(S^{d-1}\text{ 包围 }y\big)\;\Op(y)
  = e^{iq\alpha}\,\Op(y) .
  \label{eq:surround}
\end{equation}
若 $\Ssurf$ 不包围任何算符，可以收缩到一点并消失（$=1$）。

\paragraph{图像二：穿过（作用于态）。}
在正则量子化图像中，取 $\Ssurf$ 为等时面，
则 $\Uop_\alpha$ 是作用在 Hilbert 空间上的算符。
把 $\Ssurf$ 从 $\Op$ 的“过去”移动到“未来”，
就相当于让 $\Uop$ \textbf{穿过} $\Op$：
\begin{equation}
  \Uop_\alpha\;\Op(y)\;\Uop_\alpha^{-1}=e^{iq\alpha}\,\Op(y)
  \qquad\Longleftrightarrow\qquad
  \Uop_\alpha\,\Op(y)=e^{iq\alpha}\,\Op(y)\,\Uop_\alpha .
  \label{eq:passthrough}
\end{equation}

两个图像的等价性：把 \eqref{eq:passthrough} 里“穿过前”减“穿过后”，
两张等时面拼成一个包围 $\Op$ 的闭曲面，就回到 \eqref{eq:surround}。
下面这个“图”值得自己画一遍：

\begin{center}
\begin{tabular}{@{}ccc@{}}
\toprule
包围图像 & & 穿过图像 \\
\midrule
$\Ssurf$ 是包住 $\Op$ 的小球
 & $\Longleftrightarrow$
 & $\Ssurf$ 是等时面，前后各一张 \\
读出 $e^{iq\alpha}$
 & $\Longleftrightarrow$
 & 交换 $\Uop$ 与 $\Op$ 得相位 $e^{iq\alpha}$ \\
适合欧氏／路径积分
 & &
适合洛伦兹／Hilbert 空间 \\
\bottomrule
\end{tabular}
\end{center}

\begin{keypoint}
\label{kp:two-pictures}
\textbf{学会在两个图像间自由切换}是听这门课的核心技能。
课堂上说“把这条线扫过那个算符”“把这个面收缩掉”时，
用的是图像一的语言；
说“荷”“选择定则”“Hilbert 空间分解”时，用的是图像二的语言。
\end{keypoint}

\subsection{为什么“拓扑”正好对应“对称性”}

最后回答一个概念性问题：为什么拓扑性是对称性的本质，
而不只是一个附带性质？

\begin{enumerate}[label=(\arabic*)]
  \item \textbf{拓扑性 $=$ 守恒}。定理 \ref{thm:topological}：
        算符可形变 $\Leftrightarrow$ $\ext\hs j=0$。
        “荷不随时间变化”只是“可以把曲面往时间方向推”的特例。
  \item \textbf{拓扑性 $=$ 与动力学无关}。
        拓扑算符不依赖度规，因而不依赖能标；
        它在紫外与红外\textbf{同样有效}。
        这正是对称性能给出严格约束（而非近似关系）的原因，
        也是\textbf{反常匹配}（第 \ref{subsec:anomalymatch} 小节）
        能成立的根据。
  \item \textbf{拓扑性 $=$ 可以放在任何地方}。
        既然能任意形变，就能支撑在任意维数的子流形上——
        只要那个维数使得“包围／连接”有意义。
        这直接给出下一节的推广。
\end{enumerate}

% =====================================================================
\section{从 0-form 到 \texorpdfstring{$p$}{p}-form：广义全局对称性}
\label{sec:pform}

现在做推广。所有材料都已就位，推广本身只是\textbf{改一个数字}。

\subsection{推广的逻辑}

回顾普通（$0$-form）对称性的三个要素：
\begin{center}
\begin{tabular}{@{}ll@{}}
\toprule
要素 & $0$-form 情形 \\
\midrule
守恒流 & 1-形式 $j$，$\ext\hs j=0$ \\
对称算符支撑 & $(d-1)$ 维闭曲面 $\Ssurf_{d-1}$（余维 1）\\
带荷对象 & $0$ 维局域算符 $\Op(y)$ \\
连接数条件 & $(d-1)+0+1=d$ $\checkmark$ \\
\bottomrule
\end{tabular}
\end{center}

把最后一行的“$0$”换成“$q$”，其他各行随之调整，就得到定义。

\begin{definition}[$q$-form 全局对称性]
\label{def:qform}
$d$ 维 QFT 具有\textbf{$q$-form 全局对称性}（\en{$q$-form global symmetry}）
群 $G$，指存在一族算符 $\Uop_g(\Mfd_{d-q-1})$，$g\in G$，满足：
\begin{enumerate}[label=(\arabic*)]
  \item 支撑在\textbf{余维 $q+1$}（即维数 $d-q-1$）的闭子流形上；
  \item \textbf{拓扑}（定义 \ref{def:topological}）；
  \item 融合遵循群乘法：
        $\Uop_g\times \Uop_h=\Uop_{gh}$（可逆情形）；
  \item 作用在\textbf{$q$ 维带荷算符}上，通过连接数给出相位／作用。
\end{enumerate}
对连续 $G=\Uone$，存在守恒的 $(q+1)$-形式电流 $J$：
\begin{equation}
  \boxed{\ \ext\hs J=0,
  \qquad
  \Uop_\alpha(\Mfd_{d-q-1})
  =\exp\!\left[i\alpha\!\int_{\Mfd_{d-q-1}}\hs J\right]\ }
  \label{eq:qformop}
\end{equation}
其中 $\hs J\in\Omega^{d-q-1}$，次数与支撑维数匹配。$\checkmark$
\end{definition}

\begin{keypoint}
\label{kp:dimensioncheck}
\textbf{维数自查三连}。看到任何 $q$-form 对称性的陈述，
立刻检查：
\begin{align}
  \text{电流次数}&=q+1, \\
  \text{对称算符维数}&=d-q-1
  \quad(\text{余维 } q+1), \\
  \text{带荷对象维数}&=q,
  \qquad (d-q-1)+q+1=d\ \checkmark
\end{align}
三者必须同时对上。这条自查能挡掉初学阶段 90\% 的错误。
\end{keypoint}

\subsection{字典表}

\begin{table}[htbp]
\centering
\caption{$q$-form 对称性字典（$d$ 维时空）。}
\label{tab:pform}
\begin{tabular}{@{}lllll@{}}
\toprule
$q$ & 电流 & 对称算符支撑 & 带荷对象 & 典型例子 \\
\midrule
$0$ & 1-形式 $j$ & $(d-1)$ 维（余维 1） & 点（局域算符）
    & 粒子数 $\Uone$、Ising $\ZZ_2$ \\
$1$ & 2-形式 $J$ & $(d-2)$ 维（余维 2） & 线（Wilson 线）
    & Maxwell 电／磁对称性、$\SUgrp(N)$ 中心 \\
$2$ & 3-形式 $J$ & $(d-3)$ 维（余维 3） & 面（弦世界面）
    & 2-形式规范场理论、轴子弦 \\
$d-1$ & $d$-形式 & $0$ 维（点） & $(d-1)$ 维
    & “宇宙分解”，见第 \ref{sec:confusions} 节 Q\ref{q:dminus1} \\
\bottomrule
\end{tabular}
\end{table}

\begin{example}[$d=4$，$q=1$：本课程主力]
$d=4$ 中的 1-form 对称性：
电流是 2-形式，对称算符支撑在 $4-1-1=2$ 维闭曲面上，
带荷对象是 1 维的线算符。
第 \ref{subsec:maxwell} 小节将完整实现这个例子。
\end{example}

\subsection{高阶形式对称性必定是阿贝尔的}
\label{subsec:abelian}

一个漂亮而重要的结构性结果：

\begin{proposition}
\label{prop:abelian}
对 $q\ge1$，$q$-form 对称性群必为阿贝尔群。
\end{proposition}

\begin{proof}[物理论证]
对 $q\ge1$，对称算符支撑在维数 $d-q-1\le d-2$ 的子流形上，
即\textbf{余维至少 2}。余维 $\ge2$ 的两个子流形在 $d$ 维空间中
可以\textbf{互相绕过而不相交}：想象 $d=3$ 中两条直线（余维 2），
它们可以平移错开。

因此对任意两个对称算符 $\Uop_g(\Mfd)$、$\Uop_h(\Mfd')$，
总能连续地把它们交换位置而不让支撑集相交；
由拓扑性，关联函数在此过程中不变，故
\begin{equation}
  \Uop_g\,\Uop_h=\Uop_h\,\Uop_g
  \qquad\Longrightarrow\qquad
  gh=hg .
\end{equation}

对比 $q=0$：对称算符是余维 1 的（等时面），
两张等时面\textbf{无法}互相绕过——它们有确定的时间先后顺序，
所以乘法可以不交换，允许非阿贝尔群。
\end{proof}

\begin{keypoint}
命题 \ref{prop:abelian} 解释了为什么文献中 1-form 对称性总是
$\Uone$、$\ZZ_N$ 这类阿贝尔群，
而 0-form 可以是 $\SUgrp(2)$ 之类。
\textbf{“余维 1 才有时间顺序”这个直觉}值得记住，
它在讨论不可逆对称性时还会用到。
\end{keypoint}

\subsection{离散情形与 \texorpdfstring{$\ZZ_N$}{ZN}}

对离散群没有流，但定义 \ref{def:qform} 的 (1)(2)(3)(4) 仍然成立。
最常见的是 $\ZZ_N$ 的 $q$-form 对称性：
算符 $\Uop_k(\Mfd_{d-q-1})$，$k\in\ZZ_N$，满足
\begin{equation}
  \Uop_k\times \Uop_l=\Uop_{k+l\ \mathrm{mod}\ N},
  \qquad
  \big(\Uop_1\big)^N=\mathbf 1 .
\end{equation}
作用在带荷 $n$ 的 $q$ 维算符上给相位 $e^{2\pi i kn/N}$。
第 \ref{subsec:center} 小节的 $\SUgrp(N)$ 中心对称性正是此类。

\subsection{不可逆与高群：概念地图}
\label{subsec:beyond}

为免第一堂课被术语淹没，这里给一张\textbf{概念地图}，
说明各名词的相互关系。它们都是“对称性 $=$ 拓扑算符”的不同放松方向。

\begin{center}
\begin{tabular}{@{}lll@{}}
\toprule
名称 & 放松了什么 & 本讲义位置 \\
\midrule
$q$-form 对称性 & 带荷对象可以是 $q$ 维
  & 第 \ref{sec:pform} 节 \\
高群（\en{higher group}） & 不同 $q$ 的对称性非平凡地纠缠在一起
  & 本小节（仅概念）\\
不可逆对称性 & 融合规则无逆元
  & 第 \ref{subsec:ising} 小节、Q\ref{q:noninv} \\
范畴型对称性（\en{categorical}） & 用融合范畴而非群来组织全部拓扑算符
  & 第 \ref{sec:further} 节文献 \\
\bottomrule
\end{tabular}
\end{center}

\begin{remark}[高群一句话]
若把一个 0-form 对称性的背景场 $A$ 打开后，
1-form 对称性的背景场 $B$ 的规范变换里出现了 $A$
（形如 $B\to B+\ext\lambda+\text{（含 }A\text{ 的项）}$），
则两者不能独立处理，合起来叫\textbf{2-群}。
第一次听课时只需知道这个词的存在。
\end{remark}

% =====================================================================
\section{荷与缺陷}
\label{sec:charges-defects}

本节澄清一组极易混淆的概念：\textbf{对称算符}、\textbf{带荷算符}、
\textbf{缺陷}。三者维数不同、角色不同，
课堂上若混用会完全听不懂。

\subsection{三类对象的严格区分}

\begin{center}
\begin{tabular}{@{}llll@{}}
\toprule
 & 对称算符 & 带荷算符 & 缺陷 \\
 & \en{symmetry operator} & \en{charged operator} & \en{defect} \\
\midrule
维数（$q$-form） & $d-q-1$ & $q$ & 多为 $d-q-1$，但\textbf{不闭合} \\
是否拓扑 & \textbf{是} & 一般\textbf{否} & 视情形 \\
是否可形变 & 是 & 否（期望值依赖形状） & 端点／边界固定 \\
作用 & 测量／作用于他者 & 被测量 & 修改真空／边界条件 \\
Maxwell 例子 & $e^{i\alpha\oint_{\Ssurf_2}\hs F/e^2}$ & $W_n(C)$ & 开放的对称算符 \\
Ising 例子 & $\prod_j\sigma^x_j$ & $\sigma^z_j$ & 畴壁 \\
\bottomrule
\end{tabular}
\end{center}

\subsection{荷：由连接数定义}

\begin{definition}[$q$-form 荷]
设 $\Op(\Mfd_q)$ 是 $q$ 维算符，$\Uop_\alpha(\Mfd_{d-q-1})$ 是
$q$-form 对称算符。称 $\Op$ 带荷 $n$，若
\begin{equation}
  \boxed{\
  \big\langle \Uop_\alpha(\Mfd_{d-q-1})\;\Op(\Mfd_q)\;\cdots\big\rangle
  = e^{\,i\alpha\,n\,\Link(\Mfd_{d-q-1},\,\Mfd_q)}\;
  \big\langle \Op(\Mfd_q)\cdots\big\rangle \ }
  \label{eq:chargedef-q}
\end{equation}
\end{definition}

这是 \eqref{eq:surround} 的直接推广：$q=0$ 时“连接”就是“包围”。

\begin{keypoint}
\label{kp:linking-is-charge}
\textbf{荷不是“在算符上贴的标签”，而是一个连接数的读数。}
这个转变是理解广义对称性的心理门槛。
一旦接受它，就不会再问“一条线怎么可能带荷”——
因为荷本来就是“别的东西绕着它转一圈得到的相位”。
\end{keypoint}

\subsection{缺陷：把对称算符“打开”}
\label{subsec:defects}

\begin{definition}[对称缺陷]
把对称算符 $\Uop_g$ 支撑在一个\textbf{带边}的子流形上
（即不闭合），得到的对象称\textbf{对称缺陷}
（\en{symmetry defect}）或\textbf{扭曲缺陷}（\en{twist defect}）。
其\textbf{边界}是一个更低维的对象，
在该边界上理论被修改。
\end{definition}

\begin{example}[Ising 畴壁：最直观的缺陷]
\label{ex:domainwall}
1$+$1 维 Ising 链（第 \ref{subsec:ising} 小节）的 $\ZZ_2$ 对称算符是
\begin{equation}
  \Uop=\prod_{j=-\infty}^{+\infty}\sigma^x_j
  \qquad(\text{支撑在整条空间切片上，闭合}).
\end{equation}
把乘积\textbf{截断}到半无穷区间：
\begin{equation}
  \Uop_{\text{半}}=\prod_{j\le 0}\sigma^x_j
  \qquad(\text{不闭合}) .
\end{equation}
$\Uop_{\text{半}}$ \textbf{不再}与 $H$ 交换：
它在 $j=0$ 与 $j=1$ 之间的键上留下一个\textbf{畴壁}
（\en{domain wall}）。这个畴壁就是缺陷的端点，
它是一个真实的、局域的激发（有能量代价 $\sim J$）。
\end{example}

\begin{keypoint}
\label{kp:defect-vs-operator}
\textbf{“闭合 $\Rightarrow$ 对称算符（拓扑、无能量代价）；
带边 $\Rightarrow$ 缺陷（边界是物理激发）”}。
这条对应关系在所有例子中都成立，是记住区别的最好办法。
在 $q$-form 情形，$q$-form 对称缺陷的边界是一个
$(d-q-2)$ 维对象。
\end{keypoint}

\subsection{扭曲边界条件与扭曲扇区}

缺陷的另一个用途：定义\textbf{扭曲扇区}（\en{twisted sector}）。
把对称缺陷沿\textbf{时间}方向插入（即支撑在
“空间的一部分 $\times$ 全部时间”上），等价于对场施加
\textbf{扭曲边界条件}（\en{twisted boundary condition}）：
\begin{equation}
  \phi(x+L)=g\cdot\phi(x)
  \qquad(g\in G).
  \label{eq:twisted}
\end{equation}
在配分函数语言中，
\begin{equation}
  Z_g=\Tr\!\left[g\,e^{-\beta H}\right]
  \quad\text{（把 }g\text{ 插在时间圈上）},
  \label{eq:Zg}
\end{equation}
而\textbf{把 $g$ 插在空间圈上}给出扭曲扇区的 Hilbert 空间
$\Ham_g$。两者由\textbf{模变换}（\en{modular transformation}）联系
——这在 2 维 CFT 中是标准工具，也是判断不可逆对称性的实用手段。

\subsection{规范化：把全局对称性“求和掉”}
\label{subsec:gauging}

\begin{definition}[规范化]
\textbf{规范化}（\en{gauging}）一个（有限）$q$-form 对称性 $G$，
指在路径积分中\textbf{对所有可能的对称缺陷网络求和}：
\begin{equation}
  Z_{\text{gauged}}
  =\frac{1}{|G|^{\#}}\sum_{\text{网络}}\;Z\big[\text{缺陷网络}\big].
  \label{eq:gauge-sum}
\end{equation}
等价地：引入 $G$ 的 $(q+1)$-形式规范场并对其求和。
\end{definition}

规范化的三条重要后果：

\begin{enumerate}[label=(\arabic*)]
  \item \textbf{原来的带荷算符被投影掉}。求和把非中性算符的
        期望值平均为零，于是它们不再是好算符。
  \item \textbf{出现新的“对偶”对称性}。
        规范化 $d$ 维中的 $q$-form 群 $G$（阿贝尔），
        得到一个新的 $(d-q-2)$-form 对称性，群为 $\hat G$（对偶群）。
        \begin{equation}
          \text{gauge } q\text{-form } G
          \;\Longrightarrow\;
          \text{新的 } (d-q-2)\text{-form } \hat G .
          \label{eq:dualsym}
        \end{equation}
        \emph{检查 $q=0$、$d=2$、$G=\ZZ_2$}：新对称性是
        $(2-0-2)=0$-form 的 $\ZZ_2$——正是 Ising 的
        Kramers--Wannier 对偶对称性。$\checkmark$
        （第 \ref{subsec:ising} 小节）
  \item \textbf{存在反常时无法规范化}。这就是 't Hooft 反常的定义
        （第 \ref{subsec:thooft} 小节）。
\end{enumerate}

\begin{exercise}
\label{ex:gaugecheck}
用 \eqref{eq:dualsym} 验证：
(a) 在 $d=4$ 中规范化一个 0-form $\ZZ_N$，
得到什么形式的对称性？（答：$(4-0-2)=2$-form $\ZZ_N$。）\\
(b) 在 $d=4$ 中规范化一个 1-form $\ZZ_N$，得到什么？
（答：$(4-1-2)=1$-form $\ZZ_N$——\textbf{形式不变}，
这正是 4 维中 $\Uone$ 与 $\Uone/\ZZ_N$ 规范理论互相转化的原因，
见第 \ref{subsec:duality} 小节。）\\
(c) 在 $d=3$ 中规范化 0-form $\ZZ_N$，得到 $1$-form $\ZZ_N$；
说明这与 3 维 $\ZZ_N$ 规范理论（Dijkgraaf--Witten 理论）
中 Wilson 线与 't Hooft 线的辫结相位如何联系。
\end{exercise}

% =====================================================================
\section{典型例子}
\label{sec:examples}

本节是全讲义重心。三个例子分工如下：
\textbf{Ising} 演示 0-form 对称性的全部要素（含不可逆对称性的最简实例）；
\textbf{Maxwell} 给出 1-form 对称性的\textbf{完整推导}；
\textbf{$\Uone$ 对偶}提供把两者串起来的直觉。

\subsection{例一：Ising 模型与 \texorpdfstring{$\ZZ_2$}{Z2}}
\label{subsec:ising}

\subsubsection{模型与对称算符}

取 1$+$1 维横场量子 Ising 链，格点 $j$ 上是自旋 $1/2$：
\begin{equation}
  H=-J\sum_j \sigma^z_j\sigma^z_{j+1}-h\sum_j\sigma^x_j,
  \qquad J,h>0 .
  \label{eq:ising}
\end{equation}

\begin{proposition}[$\ZZ_2$ 对称性]
算符
\begin{equation}
  \boxed{\ \Uop=\prod_j\sigma^x_j\ }
  \qquad\text{满足}\qquad
  \Uop^2=1,\qquad [\Uop,H]=0 .
  \label{eq:isingU}
\end{equation}
\end{proposition}

\begin{proof}
$\Uop^2=\prod_j(\sigma^x_j)^2=1$。$\checkmark$
与 $H$ 交换：$\sigma^x_j$ 与自身交换，故第二项显然不变。
对第一项，用 $\sigma^x\sigma^z\sigma^x=-\sigma^z$ 得
\begin{equation}
  \Uop\,\sigma^z_j\sigma^z_{j+1}\,\Uop^{-1}
  =\big(-\sigma^z_j\big)\big(-\sigma^z_{j+1}\big)
  =\sigma^z_j\sigma^z_{j+1} . \checkmark
\end{equation}
\end{proof}

\paragraph{维数记账。}$d=2$（一维空间 $+$ 时间），$q=0$。
对称算符应支撑在 $d-q-1=1$ 维流形上：
在固定时刻，$\Uop$ 是\textbf{整条空间切片}上所有 $\sigma^x$ 的乘积——
确实是 1 维的。$\checkmark$（要点 \ref{kp:dimensioncheck}）
带荷对象是 0 维的局域算符：$\sigma^z_j$ 带荷 $1$（$\ZZ_2$ 奇），
因为 $\Uop\sigma^z_j\Uop^{-1}=-\sigma^z_j$。

\subsubsection{两个相与自发破缺}

\begin{itemize}
  \item \textbf{有序相}（$J\gg h$）：基态近似为
        $|\!\uparrow\uparrow\uparrow\cdots\rangle$ 与
        $|\!\downarrow\downarrow\downarrow\cdots\rangle$，
        两重退化，$\langle\sigma^z\rangle\ne0$。
        $\Uop$ 把两个基态\textbf{互换}，故
        \textbf{$\ZZ_2$ 自发破缺}（\en{spontaneously broken}）。
  \item \textbf{无序相}（$h\gg J$）：基态唯一，
        近似 $\bigotimes_j|\rightarrow\rangle_j$（$\sigma^x$ 的本征态），
        $\langle\sigma^z\rangle=0$，$\Uop|0\rangle=+|0\rangle$，
        对称性\textbf{未破缺}。
  \item 两相之间在 $J=h$ 有二阶相变，连续极限是
        $c=1/2$ 的 Ising 共形场论（\en{Ising CFT}）。
\end{itemize}

\begin{keypoint}
\label{kp:ssb-0form}
0-form 对称性自发破缺的\textbf{判据}：
存在带荷的\emph{局域}算符具有非零期望值
（\textbf{序参量}，\en{order parameter}）：
$\langle\Op\rangle\ne0$ 且 $\Op$ 带荷。
等价地：基态退化，对称算符在基态间置换。
第 \ref{subsec:confinement} 小节会看到 1-form 情形的类似判据，
只不过序参量从局域算符换成\textbf{线算符}。
\end{keypoint}

\subsubsection{缺陷：畴壁}

按第 \ref{subsec:defects} 小节，把 \eqref{eq:isingU} 截断：
\begin{equation}
  \Uop_{\le0}=\prod_{j\le 0}\sigma^x_j .
\end{equation}
它在 $j=0,1$ 之间制造一个\textbf{畴壁}。在有序相中，
畴壁是一个真实的、可局域化的激发（能量 $\sim 2J$），
它\textbf{就是缺陷的端点}。

这个例子清楚展示要点 \ref{kp:defect-vs-operator}：
\begin{itemize}
  \item 完整乘积（闭合）$\Rightarrow$ 对称算符，与 $H$ 交换，无能量代价；
  \item 截断乘积（带边）$\Rightarrow$ 缺陷，端点是有能量的激发。
\end{itemize}

\subsubsection{Kramers--Wannier 对偶与规范化}
\label{subsec:KW}

定义\textbf{对偶变量}，住在链的\textbf{键}（即 $j+\tfrac12$）上：
\begin{equation}
  \mu^x_{j+\frac12}\equiv\sigma^z_j\sigma^z_{j+1},
  \qquad
  \mu^z_{j+\frac12}\equiv\prod_{k\le j}\sigma^x_k .
  \label{eq:KWvars}
\end{equation}
可以验证 $\mu^{x,z}$ 满足同样的 Pauli 代数。代入 \eqref{eq:ising}：
\begin{equation}
  H=-J\sum \mu^x_{j+\frac12}
    -h\sum \mu^z_{j-\frac12}\mu^z_{j+\frac12} ,
  \label{eq:isingdual}
\end{equation}
即\textbf{同一个模型，但 $J\leftrightarrow h$ 互换}。

\begin{keypoint}
\label{kp:KW}
Kramers--Wannier 对偶（\en{Kramers--Wannier duality}）$J\leftrightarrow h$：
\begin{itemize}
  \item 把\textbf{有序相映到无序相}，反之亦然；
  \item 在\textbf{临界点 $J=h$ 处成为理论的自映射}——
        这是一个\emph{对称性}！
  \item 注意 \eqref{eq:KWvars} 中 $\mu^z$ 是一个\textbf{半无穷乘积}，
        即一个\textbf{缺陷}。所以对偶算符本质上是非局域的。
\end{itemize}
在现代语言中：KW 对偶 $=$ \textbf{规范化 $\ZZ_2$ 对称性}
（第 \ref{subsec:gauging} 小节）再作场重命名。
由 \eqref{eq:dualsym}，$d=2$、$q=0$、$G=\ZZ_2$ 规范化后得到
$(2-0-2)=0$-form 的 $\hat{\ZZ_2}=\ZZ_2$——即对偶模型中新的 $\ZZ_2$。
$\checkmark$
\end{keypoint}

\subsubsection{不可逆对称性：Ising 临界点的三条拓扑线}
\label{subsec:isingTDL}

在临界点（Ising CFT）上，$d=2$ 中的 0-form 对称算符支撑在 1 维流形上，
即\textbf{拓扑缺陷线}（\en{topological defect line}, TDL）。
Ising CFT 恰好有\textbf{三条}：
\begin{equation}
  \mathbf 1\ (\text{恒等}),
  \qquad
  \eta\ (\ZZ_2\ \text{对称线}),
  \qquad
  \mathcal N\ (\text{KW 对偶线}) .
\end{equation}
它们的融合规则（\eqref{eq:fusion} 的具体实现）为
\begin{equation}
  \boxed{\
  \eta\times\eta=\mathbf 1,
  \qquad
  \eta\times\mathcal N=\mathcal N\times\eta=\mathcal N,
  \qquad
  \mathcal N\times\mathcal N=\mathbf 1+\eta\ }
  \label{eq:isingfusion}
\end{equation}

\begin{keypoint}
\label{kp:noninvertible-first}
\textbf{$\mathcal N$ 没有逆元}，因此是\textbf{不可逆对称性}。
理由很直接：若存在 $\mathcal N^{-1}$，则由
$\mathcal N\times\mathcal N=\mathbf 1+\eta$ 两边左乘 $\mathcal N^{-1}$ 得
$\mathcal N=\mathcal N^{-1}+\mathcal N^{-1}\eta$，
但融合系数是\textbf{非负整数}，无法让一条线等于两条线之和。
$\mathcal N$ 的“量子维数”是 $\sqrt2$（不是整数），
这是不可逆性的标志。
\end{keypoint}

\paragraph{$\mathcal N$ 的作用与选择定则。}
Ising CFT 有三个初级场：$\mathbf 1$、能量算符 $\varepsilon$、
自旋算符 $\sigma$。三条线在它们上的作用（本征值）为
\begin{center}
\begin{tabular}{@{}lccc@{}}
\toprule
作用于 & $\mathbf 1$ & $\varepsilon$ & $\sigma$ \\
\midrule
$\eta$ & $+1$ & $+1$ & $-1$ \\
$\mathcal N$ & $+\sqrt2$ & $-\sqrt2$ & $\bm 0$ \\
\bottomrule
\end{tabular}
\end{center}
最后一格的 $0$ 是关键：\textbf{$\mathcal N$ 把 $\sigma$ 消灭}。
一个可逆算符不可能有零本征值——这又一次证明了不可逆性，
并直接给出\textbf{选择定则}：
把 $\mathcal N$ 收缩到包围 $\sigma$ 的小圈上，关联函数为零，
从而某些关联函数被迫消失。这类约束是普通群论\textbf{看不到}的，
也是不可逆对称性的实用价值。

\begin{remark}
$\eta$ 在 $\sigma$ 上取 $-1$ 说明 $\sigma$ 是 $\ZZ_2$ 奇的
（对应格点上的 $\sigma^z$），
在 $\varepsilon$ 上取 $+1$ 说明能量算符是 $\ZZ_2$ 偶的
（对应格点上的 $\sigma^x$ 或 $\sigma^z\sigma^z$）。
这与第 \ref{subsec:ising} 小节开头的格点分析完全一致。$\checkmark$
\end{remark}

\begin{exercise}
(a) 用 \eqref{eq:isingfusion} 计算 $\mathcal N\times\mathcal N\times\mathcal N$，
并验证结果与 $2\mathcal N$ 一致。\\
(b) 由表中数据验证融合规则在“本征值相乘”层面自洽：
$(\pm\sqrt2)^2=2=1+1$ 对 $\mathbf 1,\varepsilon$ 成立；
$0^2=0=1+(-1)$ 对 $\sigma$ 成立。$\checkmark$
这个自洽检查是判断融合规则记对没记对的好办法。
\end{exercise}

\subsection{例二：Maxwell 理论中的 1-form 对称性（完整推导）}
\label{subsec:maxwell}

这是本讲义最重要的推导，请逐步复核。

\subsubsection{设定与两个方程}

考虑 4 维（$d=4$）\textbf{纯} $\Uone$ 规范理论，
即 Maxwell 理论，\textbf{没有任何带电物质场}：
\begin{equation}
  S=-\frac{1}{4e^2}\int \ext^4x\;F_{\mu\nu}F^{\mu\nu}
   =-\frac{1}{2e^2}\int F\wedge\hs F ,
  \qquad F=\ext A .
  \label{eq:maxwellaction}
\end{equation}

理论满足两个方程，来源截然不同：

\begin{equation}
  \boxed{
  \begin{aligned}
    \text{运动方程：}&\quad \partial_\mu F^{\mu\nu}=0
      &&\Longleftrightarrow&&\ \ext\hs F=0 ,\\
    \text{Bianchi 恒等式：}&\quad \partial_{[\rho}F_{\mu\nu]}=0
      &&\Longleftrightarrow&&\ \ext F=0 .
  \end{aligned}}
  \label{eq:twoequations}
\end{equation}

\begin{keypoint}
\label{kp:two-sources}
\textbf{请牢记这两行的不对称性}：
\begin{itemize}
  \item $\ext\hs F=0$ 来自\textbf{变分作用量}，
        因此\textbf{加入物质场后会被修改}（源项出现）。
  \item $\ext F=0$ 来自 $F=\ext A$ 与 $\ext^2=0$，
        是\textbf{恒等式}，\textbf{不依赖任何动力学}，
        只有在允许磁单极子（$A$ 不整体存在）时才被修改。
\end{itemize}
两个 1-form 对称性分别源于这两行，
命运也因此不同（第 \ref{subsec:matter} 小节）。
\end{keypoint}

\subsubsection{两个守恒的 2-形式电流}

按定义 \ref{def:qform}，$d=4$、$q=1$ 需要 2-形式电流。
\eqref{eq:twoequations} 恰好提供两个：

\begin{proposition}[电与磁 1-form 对称性]
定义
\begin{equation}
  \hs J_e\equiv\frac{1}{e^2}\,\hs F
  \qquad\text{与}\qquad
  \hs J_m\equiv\frac{1}{2\pi}\,F ,
  \label{eq:twocurrents}
\end{equation}
两者都是 2-形式且都闭：
\begin{equation}
  \ext\hs J_e=\frac{1}{e^2}\ext\hs F=0\ (\text{运动方程}),
  \qquad
  \ext\hs J_m=\frac{1}{2\pi}\ext F=0\ (\text{Bianchi}) .
\end{equation}
\end{proposition}

\paragraph{维数自查（要点 \ref{kp:dimensioncheck}）。}
$q=1$，$d=4$：电流是 $q+1=2$-形式 $\checkmark$；
$\hs J$ 是 $d-q-1=2$-形式 $\checkmark$（上面两个都是 2-形式）；
对称算符支撑在 2 维闭曲面上 $\checkmark$；
带荷对象是 1 维的线 $\checkmark$；
连接数 $2+1+1=4$ $\checkmark$。全部对上。

\begin{definition}[Maxwell 的两个对称算符]
对任意 2 维闭曲面 $\Ssurf_2\subset\Mfd_4$：
\begin{equation}
  \boxed{\
  \Uop_e^{(\alpha)}(\Ssurf_2)
  =\exp\!\left[\frac{i\alpha}{e^2}\oint_{\Ssurf_2}\hs F\right],
  \qquad
  \Uop_m^{(\beta)}(\Ssurf_2)
  =\exp\!\left[\frac{i\beta}{2\pi}\oint_{\Ssurf_2}F\right] \ }
  \label{eq:maxwellops}
\end{equation}
由定理 \ref{thm:topological}，两者都是\textbf{拓扑算符}。
\end{definition}

\subsubsection{推导：电对称性作用在 Wilson 线上}
\label{subsec:wilson-charge}

现在完整推导 $\Uop_e$ 在 Wilson 线上的作用。

\paragraph{第 1 步：把 Wilson 线插入作用量。}
插入 $W_n(C)=\exp[in\oint_C A]$ 等价于在作用量中加一项
$n\oint_C A$。用 Poincaré 对偶 \eqref{eq:poincare} 把线积分写成时空积分：
\begin{equation}
  n\oint_C A=n\int_{\Mfd_4}\delta_C\wedge A ,
  \label{eq:wilsonsource}
\end{equation}
其中 $\delta_C$ 是支撑在 $C$ 上的 \textbf{3-形式}
（$d-k=4-1=3$ $\checkmark$）。

\paragraph{第 2 步：变分。}对 $A\to A+\delta A$ 变分总作用量
$S_{\rm tot}=-\frac{1}{2e^2}\int F\wedge\hs F+n\int\delta_C\wedge A$。
动能项给出 $\propto\frac{1}{e^2}\ext\hs F$，源项给出 $n\,\delta_C$，
于是运动方程被修改为
\begin{equation}
  \boxed{\ \frac{1}{e^2}\,\ext\hs F=n\,\delta_C\ }
  \label{eq:modifiedEOM}
\end{equation}
（整体符号取决于定向与号差约定；下面只用它的\emph{绝对}效果。）

\paragraph{第 3 步：用 Stokes 定理算荷。}
取 $\Ssurf_2=\partial V_3$，其中 $V_3$ 是一个 3 维区域，
与 $C$ \textbf{恰好相交一次}（即 $\Ssurf_2$ 与 $C$ 的连接数为 1）。则
\begin{align}
  \frac{1}{e^2}\oint_{\Ssurf_2}\hs F
  &\overset{\text{Stokes}}{=}
   \frac{1}{e^2}\int_{V_3}\ext\hs F
  \notag\\
  &\overset{\eqref{eq:modifiedEOM}}{=}
   n\int_{V_3}\delta_C
  \notag\\
  &\overset{\text{Poincaré 对偶}}{=}
   n\cdot\#\big(V_3\cap C\big)
  \;=\;n .
  \label{eq:fluxisn}
\end{align}

\paragraph{第 4 步：结论。}代入 \eqref{eq:maxwellops}：
\begin{equation}
  \boxed{\
  \Uop_e^{(\alpha)}(\Ssurf_2)\;W_n(C)
  = e^{\,i\alpha n\,\Link(\Ssurf_2,C)}\;W_n(C)\ }
  \label{eq:UeonW}
\end{equation}
这正是 \eqref{eq:chargedef-q} 的形式：
\textbf{Wilson 线 $W_n$ 在电 1-form 对称性下带荷 $n$}。$\blacksquare$

\begin{remark}[Gauss 律图像]
\label{rmk:gauss}
\eqref{eq:fluxisn} 的物理内容就是高中电磁学的 Gauss 律。
取 $\Ssurf_2$ 为固定时刻包围电荷的球面 $S^2$，
则 $\frac{1}{e^2}\oint_{S^2}\hs F$ 恰为\textbf{电通量}
$\oint_{S^2}\bm E\cdot \ext\bm S$，而 Gauss 律说它等于所包围的电荷 $n$。

于是 \eqref{eq:UeonW} 的意思是：
\begin{center}
\emph{“电 1-form 对称算符”就是“测量电通量的算符”，
而它的荷就是被套住的电荷。}
\end{center}
\textbf{这个例子说明广义对称性并不神秘}：
它把一个人人都会的操作（用高斯面测电荷）
重新识别为一个对称性的作用。新的地方在于：
把 $S^2$ 换成任意 2 维闭曲面、把“包围”换成“连接”，
并认识到 $W_n(C)$ 是被作用的\emph{线}算符而不是点算符。
\end{remark}

\subsubsection{磁对称性与 't Hooft 线}
\label{subsec:tHooft-line}

对磁对称性，带荷对象是 \textbf{'t Hooft 线} $T_m(C)$（例 \ref{ex:thooft}）。
按定义 \eqref{eq:thooftdef}，环绕 $C$ 的小球面上磁通为 $m$：
\begin{equation}
  \frac{1}{2\pi}\oint_{\Ssurf_2}F=m\qquad
  \big(\Link(\Ssurf_2,C)=1\big),
\end{equation}
于是直接由 \eqref{eq:maxwellops}
\begin{equation}
  \boxed{\
  \Uop_m^{(\beta)}(\Ssurf_2)\;T_m(C)
  =e^{\,i\beta m\,\Link(\Ssurf_2,C)}\;T_m(C)\ }
  \label{eq:UmonT}
\end{equation}

注意这里\textbf{不需要}修改运动方程：$m$ 的整数性来自
第 \ref{subsec:forms} 小节 $S^2$ 上 $H^2\ne0$ 的拓扑论证
（Dirac 量子化），而非来自动力学。
这与要点 \ref{kp:two-sources} 完全吻合。$\checkmark$

\begin{remark}[一个值得向老师提问的细节]
$\frac{1}{2\pi}\oint_{\Ssurf_2}F\in\ZZ$ 对任意闭 $\Ssurf_2$ 成立
（$A$ 是 $\Uone$ 联络），因此 $\beta$ 是 $2\pi$ 周期的，
磁对称群确实是 $\Uone$。
但 $\frac{1}{e^2}\oint_{\Ssurf_2}\hs F$ 在\emph{纯} Maxwell 理论中
是一个\textbf{连续}的实数，其“整数性”只在 $\Ssurf_2$ 套住
Wilson 线时由 \eqref{eq:fluxisn} 保证。
因此“电 1-form 对称群是 $\Uone$ 还是 $\RR$”依赖于
如何精确设定带荷线算符谱与全局结构。
文献中通常写 $\Uone_e^{(1)}\times\Uone_m^{(1)}$；
这是一个值得课上确认约定的地方。
\end{remark}

\subsubsection{加入物质场：对称性如何被破坏}
\label{subsec:matter}

现在检验要点 \ref{kp:two-sources}。加入一个电荷为 $q$ 的
\textbf{动力学}物质场 $\psi$，则运动方程变为
\begin{equation}
  \frac{1}{e^2}\ext\hs F=\hs j_\psi\ne0
  \qquad\Longrightarrow\qquad
  \ext\hs J_e\ne0 .
  \label{eq:brokenEOM}
\end{equation}
$\hs J_e$ 不再闭 $\Rightarrow$ $\Uop_e(\Ssurf_2)$ \textbf{不再是拓扑算符}
$\Rightarrow$ 连续的电 1-form 对称性\textbf{被破坏}。

\paragraph{但残留一个 $\ZZ_q$。}
物质场电荷为 $q$，因此它能\textbf{屏蔽}（例 \ref{ex:screening}）
电荷为 $q$ 整数倍的 Wilson 线：
\begin{equation}
  W_n\ \sim\ W_{n+q}
  \qquad\Longrightarrow\qquad
  \text{线算符按}\ n\ \mathrm{mod}\ q\ \text{分类}.
\end{equation}
于是剩下的 1-form 对称群是
\begin{equation}
  \boxed{\ \Uone_e^{(1)}\ \longrightarrow\ \ZZ_q^{(1)}\ }
  \label{eq:UtoZq}
\end{equation}
特别地，若存在电荷 $1$ 的物质场（如普通 QED 中的电子，$q=1$），
则 $\ZZ_1=$ 平凡群，\textbf{电 1-form 对称性被完全破坏}。

\paragraph{磁对称性则依然存活。}$\ext F=0$ 是恒等式，
加入\emph{电}荷物质不影响它。只有引入\textbf{动力学磁单极子}
才会破坏磁 1-form 对称性。

\begin{keypoint}
\label{kp:whypure}
这解释了一个初学者常有的疑问：
\textbf{“为什么教科书里从没提过 1-form 对称性？”}
因为现实世界的 QED 有电荷 $1$ 的电子，
电 1-form 对称性被完全破坏；
而磁对称性虽然存活，却因为没有动力学磁单极子而很少被用到。
1-form 对称性在\textbf{纯规范理论}
（如格点 Yang--Mills、拓扑物态）中才显出威力。
\end{keypoint}

\subsubsection{\texorpdfstring{$\SUgrp(N)$}{SU(N)} 的 \texorpdfstring{$\ZZ_N$}{ZN} 中心对称性}
\label{subsec:center}

把上面的逻辑搬到非阿贝尔理论。取 4 维纯 $\SUgrp(N)$ Yang--Mills
（只有胶子，无基本表示物质）。

由 \eqref{eq:center}，$Z(\SUgrp(N))=\ZZ_N$。定义对称算符
$\Uop_k(\Ssurf_2)$，$k\in\ZZ_N$，作用在表示 $R$ 的 Wilson 线上：
\begin{equation}
  \boxed{\
  \Uop_k(\Ssurf_2)\;W_R(C)
  =\exp\!\left[\frac{2\pi i\,k\,n_R}{N}\Link(\Ssurf_2,C)\right] W_R(C)\ }
  \label{eq:centeraction}
\end{equation}
其中 $n_R$ 是表示 $R$ 的 $N$ 度（例 \ref{ex:Nality}）。

\begin{itemize}
  \item 基本表示：$n_R=1$，荷非零 $\Rightarrow$ 是好的序参量。
  \item 伴随表示：$n_R=0$，荷为零。这与物理一致：
        伴随 Wilson 线可以被胶子屏蔽，不携带守恒荷。
  \item 这就是\textbf{$\ZZ_N$ 1-form 中心对称性}
        （\en{$\ZZ_N$ one-form center symmetry}）。
        加入基本表示的动力学夸克（$n_R=1$）会像
        \eqref{eq:UtoZq} 一样把它完全破坏——
        这正是为什么真实 QCD 没有精确的中心对称性。
\end{itemize}

\subsubsection{禁闭 \texorpdfstring{$=$}{=} 1-form 对称性未破缺}
\label{subsec:confinement}

现在是这套语言最漂亮的物理回报。

按要点 \ref{kp:ssb-0form}，判断对称性是否自发破缺要看
\textbf{带荷算符的期望值}。对 1-form 对称性，
带荷算符是 Wilson 线，于是序参量是 $\langle W(C)\rangle$。
取 $C$ 为边长 $L$、$T$ 的大矩形圈，两种可能行为：

\begin{equation}
  \big\langle W(C)\big\rangle\ \sim\
  \begin{cases}
    e^{-\sigma\,\mathrm{Area}(C)}
      & \text{面积律（\en{area law}）}\\[4pt]
    e^{-c\,\mathrm{Perimeter}(C)}
      & \text{周长律（\en{perimeter law}）}
  \end{cases}
  \label{eq:arealaw}
\end{equation}

\begin{keypoint}
\label{kp:confinement}
\begin{center}
\begin{tabular}{@{}lll@{}}
\toprule
$\langle W(C)\rangle$ & 1-form 对称性 & 物理相 \\
\midrule
面积律 & \textbf{未}自发破缺 & \textbf{禁闭}（\en{confinement}）\\
周长律 & \textbf{已}自发破缺 & 解禁闭 / Coulomb 相 \\
\bottomrule
\end{tabular}
\end{center}
理由（直觉版）：周长律意味着可以通过重整化
（吸收掉与 $C$ 长度成正比的自能）使 $\langle W\rangle$ 保持有限非零，
即“带荷算符有非零期望值”$\Rightarrow$ 破缺。
面积律使 $\langle W\rangle$ 随面积指数衰减到零，
无法用局域重整化救回来 $\Rightarrow$ 未破缺。
\end{keypoint}

\begin{remark}[为什么这是个进步]
“禁闭”原本是个现象学描述（夸克拔不出来、有线性势 $V(r)=\sigma r$）。
现在它变成了一个\textbf{对称性破缺的判断}，
从而可以套用所有关于对称性破缺的一般理论
（序参量、Landau 论证、反常匹配）。
特别地，在纯 Yang--Mills 中 $\ZZ_N$ 是\textbf{精确}对称性，
所以“禁闭相”是一个有严格定义的相，而不只是强耦合的定性说法。
\end{remark}

\begin{remark}[光子是 1-form 对称性破缺的 Goldstone 玻色子]
在自由 Maxwell（Coulomb 相），$\langle W(C)\rangle$ 服从周长律，
故\textbf{电与磁 1-form 对称性都自发破缺}。
按与 Goldstone 定理平行的逻辑，破缺应伴随无质量激发——
那正是\textbf{光子}。
所以“为什么光子无质量”有一个对称性解释：
它是 1-form 对称性自发破缺的 Goldstone 模。
第 \ref{sec:ssb} 节会再谈。
\end{remark}

\begin{exercise}
\label{ex:maxwell}
(a) 把 \eqref{eq:fluxisn} 的每一步的形式次数写出来，
确认所有积分都在正确维数的流形上进行。\\
(b) 说明为什么 $\Uop_e$ 与 $\Uop_m$ 互相交换
（提示：命题 \ref{prop:abelian}，两个 2 维面在 4 维中可以错开）。\\
(c) 若在 $\SUgrp(N)$ 理论中加入 $N$ 度为 $2$ 的物质场，
残留的中心对称性是什么？
（答：$\ZZ_{\gcd(2,N)}$；例如 $N=4$ 时残留 $\ZZ_2$。）
\end{exercise}

\subsection{例三：\texorpdfstring{$\Uone$}{U(1)} 对偶的直觉}
\label{subsec:duality}

\subsubsection{4 维电磁对偶}

Maxwell 方程 \eqref{eq:twoequations} 在替换
\begin{equation}
  F\ \longrightarrow\ \hs F,
  \qquad \hs F\ \longrightarrow\ -F
  \label{eq:EMduality}
\end{equation}
下\textbf{互换}：运动方程 $\ext\hs F=0$ 变成 Bianchi $\ext F=0$，
反之亦然。这就是\textbf{电磁对偶}（\en{electromagnetic duality}）。

\begin{keypoint}
\label{kp:dualityswaps}
在广义对称性语言中，对偶的内容极其干净：
\begin{center}
\begin{tabular}{@{}lcl@{}}
\toprule
原理论 & $\longleftrightarrow$ & 对偶理论 \\
\midrule
运动方程 $\ext\hs F=0$ & $\longleftrightarrow$ & Bianchi $\ext F=0$ \\
电 1-form 对称性 $\Uone_e^{(1)}$ & $\longleftrightarrow$
  & 磁 1-form 对称性 $\Uone_m^{(1)}$ \\
Wilson 线 $W_n$ & $\longleftrightarrow$ & 't Hooft 线 $T_n$ \\
耦合 $e$ & $\longleftrightarrow$ & $\sim 1/e$ \\
\bottomrule
\end{tabular}
\end{center}
引入复化耦合 $\tau=\dfrac{\theta}{2\pi}+\dfrac{2\pi i}{e^2}$，
对偶变换是 $\tau\to-1/\tau$（$S$ 对偶）。
\end{keypoint}

\subsubsection{更简单的类比：2 维紧致标量与 T 对偶}
\label{subsec:compactscalar}

4 维电磁对偶的“最小模型”是 2 维紧致标量场
$\varphi\sim\varphi+2\pi$，作用量（示意）
\begin{equation}
  S=\frac{R^2}{4\pi}\int \ext\varphi\wedge\hs \ext\varphi .
\end{equation}
它有\textbf{两个} 0-form $\Uone$ 对称性，来源与 Maxwell 完全平行：

\begin{center}
\begin{tabular}{@{}lll@{}}
\toprule
 & 电流 & 守恒的原因 \\
\midrule
\textbf{平移}（动量）对称性 $\varphi\to\varphi+c$
  & $j_{\rm s}=\dfrac{R^2}{2\pi}\hs \ext\varphi$
  & 运动方程 $\ext\hs \ext\varphi=0$ \\[4pt]
\textbf{缠绕}（winding）对称性
  & $j_{\rm w}=\dfrac{1}{2\pi}\ext\varphi$
  & 恒等式 $\ext\ext\varphi=0$ \\
\bottomrule
\end{tabular}
\end{center}

\begin{keypoint}
\label{kp:parallel}
\textbf{请对比这张表与要点 \ref{kp:two-sources}}：
结构\textbf{一模一样}——
一个对称性来自\emph{运动方程}，另一个来自\emph{恒等式}。
这不是巧合，而是“对偶成对出现的对称性”的普遍模式。
掌握这个模式后，看到任何理论都可以问：
\emph{它的运动方程和它的恒等式各给我一个什么对称性？}
\end{keypoint}

\paragraph{荷与带荷算符。}
\begin{itemize}
  \item 动量对称性的荷是\textbf{动量} $n$，
        带荷算符是顶点算符 $e^{in\varphi}$
        （在 $\varphi\to\varphi+c$ 下得相位 $e^{inc}$）。
  \item 缠绕对称性的荷是\textbf{缠绕数}
        $w=\frac{1}{2\pi}\oint \ext\varphi\in\ZZ$，
        带荷算符是涡旋／无序算符。
\end{itemize}

\paragraph{T 对偶。}变换
\begin{equation}
  R\ \longrightarrow\ \frac{1}{R},
  \qquad
  n\ \longleftrightarrow\ w
  \label{eq:Tduality}
\end{equation}
是理论的同构，\textbf{交换两个 $\Uone$ 对称性}。
这与 \eqref{eq:EMduality} 交换 $\Uone_e^{(1)}\leftrightarrow\Uone_m^{(1)}$
是同一个故事的低维版本。

\subsubsection{全局结构：同一个局域理论，不同的线算符谱}

最后一个概念上重要的点。考虑 4 维规范理论，
规范群取 $\Uone$ 还是 $\Uone/\ZZ_N$？
两者的\textbf{局域}动力学（Feynman 规则、微扰展开、
$\langle F F\rangle$ 关联函数）\textbf{完全相同}，
差别只在\textbf{允许存在哪些线算符}：

\begin{itemize}
  \item 取 $\Uone$：Wilson 线 $W_n$，$n\in\ZZ$；
  \item 取 $\Uone/\ZZ_N$：允许分数电荷的线，
        但 't Hooft 线的谱相应受限。
\end{itemize}

由练习 \ref{ex:gaugecheck}(b)，在 $d=4$ 中规范化一个 1-form $\ZZ_N$
又得到一个 1-form $\ZZ_N$，所以这些不同的“全局结构”
（\en{global structure}）由\textbf{规范化 1-form 对称性的不同子群}
互相联系。

\begin{keypoint}
\label{kp:globalstructure}
\textbf{“理论”不只由拉氏量决定，还由算符谱决定。}
这是 GKSW 论文最有影响的信息之一：
一个拉氏量可能对应\textbf{若干个}不同的量子理论，
它们的区别由 1-form 对称性及其规范化选择刻画。
在 $\SUgrp(N)$ 与 $\SUgrp(N)/\ZZ_N=\mathrm{PSU}(N)$ Yang--Mills 之间、
在 $\mathcal N=4$ 超对称理论的不同全局形式之间，
这个区别都有真实的物理后果。
\end{keypoint}

% =====================================================================
\section{自发破缺与 't Hooft 反常：课前级直觉}
\label{sec:ssb}

本节目标不是把反常算清楚（那需要好几周），
而是让你在第一堂课听到 “anomaly” 时知道\textbf{它在问什么问题}。

\subsection{自发破缺的统一判据}
\label{subsec:ssb-unified}

\begin{definition}[广义自发破缺]
一个 $q$-form 对称性称为\textbf{自发破缺}的
（\en{spontaneously broken}），若存在带荷的 $q$ 维算符
$\Op(\Mfd_q)$，其期望值在“去掉与 $\Mfd_q$ 体积成正比的
局域重整化因子”后仍非零：
\begin{equation}
  \big\langle \Op(\Mfd_q)\big\rangle
  \ \sim\ e^{-c\cdot\mathrm{Vol}(\partial\Mfd_q)}
  \times(\text{非零})
  \qquad\text{（“周长律”）} .
  \label{eq:ssb-general}
\end{equation}
若衰减比这更快（如“面积律”$e^{-\sigma\mathrm{Vol}(\Mfd_q)}$），
则\textbf{未破缺}。
\end{definition}

检查两个特例：

\begin{center}
\begin{tabular}{@{}llll@{}}
\toprule
$q$ & 带荷对象 & 破缺判据 & 无质量激发（“Goldstone”） \\
\midrule
$0$ & 局域算符 $\Op(x)$ & $\langle\Op\rangle\ne0$
    & Goldstone 标量玻色子 \\
$1$ & Wilson 线 $W(C)$ & 周长律
    & 无质量规范玻色子（光子） \\
\bottomrule
\end{tabular}
\end{center}

\begin{remark}[$q=0$ 时 \eqref{eq:ssb-general} 退化为老判据]
$q=0$ 时 $\Mfd_0$ 是一个点，$\partial\Mfd_0=\emptyset$，
“体积”为零，于是 \eqref{eq:ssb-general} 就是 $\langle\Op\rangle\ne0$。
$\checkmark$ 老判据是新判据的特例。
\end{remark}

\begin{keypoint}
\label{kp:goldstone-general}
广义 Goldstone 定理的一句话版本：
\begin{center}
\emph{$q$-form 对称性自发破缺 $\Rightarrow$ 存在无质量的
$q$-form 规范场（其“Goldstone 模”）。}
\end{center}
$q=0$ 给出普通的无质量标量；
$q=1$ 给出无质量的 1-form 规范场，即\textbf{光子}。
这就是第 \ref{subsec:confinement} 小节末尾那句
“光子是 1-form Goldstone 玻色子”的来源。
\end{keypoint}

\subsection{什么是 't Hooft 反常}
\label{subsec:thooft}

\subsubsection{先分清两种“反常”}

\begin{warning}
\label{warn:twoanomalies}
“反常”这个词在文献里指两件\textbf{不同}的事，务必区分：
\begin{enumerate}[label=(\arabic*)]
  \item \textbf{ABJ 型／真反常}：对称性在量子层面
        \emph{根本不存在}，守恒律被破坏
        （如手征对称性因三角图而 $\partial_\mu j^\mu_A\propto F\tilde F$）。
        这种“对称性”不能用来做选择定则。
  \item \textbf{'t Hooft 反常（\en{'t Hooft anomaly}）}：
        对称性\emph{完全成立}，守恒律没问题；
        但它\textbf{不能被规范化}（\en{cannot be gauged}）。
        这类反常\textbf{非常有用}，因为它给出对低能物理的严格约束。
\end{enumerate}
本课程说 “anomaly” 时，绝大多数指第 (2) 种。
\end{warning}

\subsubsection{定义：背景场语言}

规范化一个对称性的第一步是\textbf{打开背景场}
（\en{background field}）。对 $q$-form 对称性 $G$，
背景场是一个 $(q+1)$-形式规范场 $B$：
\begin{equation}
  q=0:\ \text{1-形式 }A;\qquad
  q=1:\ \text{2-形式 }B;\qquad
  \text{一般 } q:\ (q+1)\text{-形式} .
\end{equation}
配分函数变成 $Z[B]$。

\begin{definition}['t Hooft 反常]
\label{def:thooft}
若不存在任何\textbf{局域反项}（\en{local counterterm}）使
$Z[B]$ 在背景规范变换 $B\to B+\ext\lambda$ 下不变，
即总有
\begin{equation}
  Z[B+\ext\lambda]=e^{\,i\,\mathcal A[B,\lambda]}\;Z[B],
  \qquad \mathcal A\ne0\ (\text{模局域反项}),
  \label{eq:anomalyphase}
\end{equation}
则称该对称性有 \textbf{'t Hooft 反常}。
\end{definition}

\begin{remark}[反常流入与 SPT 项]
\eqref{eq:anomalyphase} 的相位可以被“吸收”到
高一维的\textbf{拓扑项}中：存在 $(d+1)$ 维的
SPT／Chern--Simons 型作用量 $S_{d+1}[B]$，使得
\begin{equation}
  Z_d[B]\times e^{iS_{d+1}[B]}
\end{equation}
整体规范不变。这叫\textbf{反常流入}（\en{anomaly inflow}）。
因此“反常”被 $(d+1)$ 维拓扑项完全分类——
这是现代处理反常的标准方式。
\end{remark}

\subsubsection{最重要的推论：反常匹配}
\label{subsec:anomalymatch}

\begin{theorem}['t Hooft 反常匹配（陈述）]
\label{thm:matching}
't Hooft 反常是重整化群流的不变量：
紫外理论与红外理论的反常必须相等。
\end{theorem}

\begin{proof}[直觉论证]
反常由拓扑项 $S_{d+1}[B]$ 刻画，而拓扑项\textbf{不依赖度规}，
因而不依赖能标（第 \ref{sec:topological} 节第 (2) 条）。
系数是量子化的整数，不能连续变化。
既然它在 RG 流下不能变，紫外与红外必须给出同一个值。
\end{proof}

\begin{keypoint}
\label{kp:matching-power}
反常匹配的\textbf{威力}：它把“低能会发生什么”从猜测变成了推理。
若紫外理论有非零反常，则红外\textbf{不可能}是“平凡的有能隙唯一真空”，
必须满足以下之一：
\begin{enumerate}[label=(\arabic*)]
  \item 对称性\textbf{自发破缺}（则 Goldstone 模承载反常）；
  \item 理论在低能\textbf{无能隙}（如 CFT）；
  \item 低能是一个\textbf{拓扑场论}（TQFT），
        其反常与紫外匹配。
\end{enumerate}
\textbf{广义对称性带来的新东西}就是：把这套论证用到
高阶形式对称性上，从而对禁闭、$\theta$ 依赖、
相图结构得到全新的约束。
\end{keypoint}

\subsection{三个层次的反常例子}

\subsubsection{最初等的例子：Dirac 量子化就是一种“互不相容”}

考虑 4 维 Maxwell 的 Wilson 线 $W_n$ 与 't Hooft 线 $T_m$。
把 $W_n$ 绕着 $T_m$ 走一圈（其世界线扫出一个包围 $T_m$ 的曲面），
得到 Aharonov--Bohm 型相位
\begin{equation}
  e^{2\pi i\,n m} .
  \label{eq:diracphase}
\end{equation}
要求单值性给出 \textbf{Dirac 量子化条件} $nm\in\ZZ$。

\begin{keypoint}
\label{kp:mutual}
\eqref{eq:diracphase} 的含义是：$W$ 与 $T$ \textbf{互不局域}
（\en{mutually non-local}）。你不能同时把两者都当作“平凡的”算符。
用对称性语言：\textbf{电与磁 1-form 对称性不能同时被规范化}——
这就是它们之间存在\textbf{混合 't Hooft 反常}
（\en{mixed 't Hooft anomaly}）的最直观表现。
\end{keypoint}

\subsubsection{4 维 Maxwell 的混合反常}

把电与磁 1-form 对称性的背景场记为 2-形式 $B_e$、$B_m$。
两者的混合反常由 5 维拓扑项刻画，其结构为
\begin{equation}
  S_5\ \propto\ \int_{\Mfd_5} B_e\wedge \ext B_m
  \label{eq:5danomaly}
\end{equation}
（具体的 $2\pi$ 归一化因作者而异，请以课上约定为准）。

\paragraph{物理翻译。}\eqref{eq:5danomaly} 非零意味着：
\begin{itemize}
  \item 不能同时规范化 $\Uone_e^{(1)}$ 与 $\Uone_m^{(1)}$；
  \item 因此 Maxwell 理论\textbf{不可能}在保持两个对称性的前提下
        变成平凡的有能隙理论；
  \item 由要点 \ref{kp:matching-power}，只剩“对称性破缺”或“无能隙”
        两条路。而我们知道答案是：两者都破缺，且理论无能隙——
        \textbf{无质量光子}正是这个反常的必然后果。$\checkmark$
\end{itemize}

\begin{keypoint}
\label{kp:photon-explained}
把上面几条串起来，得到一个漂亮的论证链：
\begin{center}
混合反常 \eqref{eq:5danomaly} $\Rightarrow$
红外不能平凡 $\Rightarrow$
1-form 对称性破缺（周长律）$\Rightarrow$
存在 Goldstone 模 $\Rightarrow$
\textbf{光子无质量}。
\end{center}
这是本讲义\textbf{最值得带进课堂的一条论证}：
它把六七个概念（反常、匹配、破缺、序参量、Goldstone、拓扑算符）
串成一条可复述的因果链。
\end{keypoint}

\subsubsection{\texorpdfstring{$\theta$}{theta} 角与时间反演（仅指路）}

一个影响很大的应用：4 维 $\SUgrp(N)$ Yang--Mills 的
$\ZZ_N$ 1-form 中心对称性与时间反演／$\theta$ 周期性之间
存在混合反常。结论是：在 $\theta=\pi$ 处，
理论\textbf{不能}既保持时间反演又保持中心对称性并且有唯一有能隙真空，
从而推出 $\theta=\pi$ 处必有相变或自发破缺。
这属于课程后半段内容，第 \ref{sec:further} 节给出文献。

\subsection{小结：新旧对照}

\begin{center}
\begin{tabular}{@{}lll@{}}
\toprule
概念 & 0-form（旧） & $q$-form（新） \\
\midrule
序参量 & 局域算符期望值 & $q$ 维算符的周长律 \\
Goldstone 模 & 无质量标量 & 无质量 $q$-form 规范场 \\
反常 & $(d+1)$ 维 SPT 项 & 同样，但背景场是 $(q+1)$-形式 \\
反常匹配 & 约束低能谱 & 约束禁闭／$\theta$ 依赖／相图 \\
规范化 & 得对偶 $(d-2)$-form & 得对偶 $(d-q-2)$-form \\
\bottomrule
\end{tabular}
\end{center}

% =====================================================================
\section{课堂公式与英汉术语表}
\label{sec:glossary}

\subsection{核心公式速查}

把全篇最该带进课堂的公式集中如下。

\paragraph{（一）守恒与拓扑。}
\begin{align}
  \partial_\mu j^\mu=0
  &\ \Longleftrightarrow\ \ext\hs j=0
  && \text{（\eqref{eq:conservation}）}\\
  Q(\Ssurf_{d-1})=\int_{\Ssurf_{d-1}}\hs j
  &\ \text{只依赖 }[\Ssurf_{d-1}]\in H_{d-1}
  && \text{（定理 \ref{thm:topological}）}\\
  \Uop_\alpha(\Ssurf)&=\exp\!\Big[i\alpha\!\int_\Ssurf\hs j\Big]
  && \text{（\eqref{eq:Qsurface}）}
\end{align}

\paragraph{（二）Noether。}
\begin{align}
  j^\mu&=\PD{\Lag}{(\partial_\mu\phi^a)}\Delta^a-K^\mu
  && \text{（\eqref{eq:noethercurrent}）}\\
  j^\mu_{\Uone}&=i\big(\phi^*\partial^\mu\phi-\phi\,\partial^\mu\phi^*\big)
  && \text{（\eqref{eq:u1current}）}\\
  T^\mu{}_\nu&=\PD{\Lag}{(\partial_\mu\phi^a)}\partial_\nu\phi^a
    -\delta^\mu_\nu\Lag
  && \text{（\eqref{eq:stresstensor}）}\\
  \partial_\mu\langle j^\mu(x)\Op(y)\rangle
  &=q\,\delta^d(x-y)\langle\Op(y)\rangle
  && \text{（\eqref{eq:ward}）}
\end{align}

\paragraph{（三）$q$-form 对称性。}
\begin{align}
  \ext\hs J&=0,\qquad J\in\Omega^{q+1},\quad \hs J\in\Omega^{d-q-1}
  && \text{（\eqref{eq:qformop}）}\\
  \Uop_\alpha(\Mfd_{d-q-1})
  &=\exp\!\Big[i\alpha\!\int_{\Mfd_{d-q-1}}\hs J\Big]
  && \text{（\eqref{eq:qformop}）}\\
  \langle \Uop_\alpha(\Mfd_{d-q-1})\Op(\Mfd_q)\rangle
  &=e^{i\alpha n\Link}\langle\Op(\Mfd_q)\rangle
  && \text{（\eqref{eq:chargedef-q}）}\\
  \dim\Mfd_{\rm sym}+\dim\Mfd_{\rm charge}+1&=d
  && \text{（\eqref{eq:linkdim}）}
\end{align}

\paragraph{（四）Maxwell 的两个对称性。}
\begin{align}
  \Uop_e^{(\alpha)}(\Ssurf_2)
  &=\exp\!\Big[\tfrac{i\alpha}{e^2}\oint_{\Ssurf_2}\hs F\Big],
  & \Uop_e\,W_n&=e^{i\alpha n\Link}W_n
  && \text{（\eqref{eq:UeonW}）}\\
  \Uop_m^{(\beta)}(\Ssurf_2)
  &=\exp\!\Big[\tfrac{i\beta}{2\pi}\oint_{\Ssurf_2}F\Big],
  & \Uop_m\,T_m&=e^{i\beta m\Link}T_m
  && \text{（\eqref{eq:UmonT}）}
\end{align}
加入电荷 $q$ 的动力学物质：$\Uone_e^{(1)}\to\ZZ_q^{(1)}$
（\eqref{eq:UtoZq}）。

\paragraph{（五）Ising 的融合规则。}
\begin{equation}
  \eta^2=\mathbf 1,\qquad
  \eta\,\mathcal N=\mathcal N,\qquad
  \mathcal N^2=\mathbf 1+\eta
  \qquad \text{（\eqref{eq:isingfusion}）}
\end{equation}

\paragraph{（六）规范化与对偶。}
\begin{equation}
  \text{gauge }q\text{-form }G\ \text{in }d\text{ dims}
  \ \Longrightarrow\
  (d-q-2)\text{-form }\hat G
  \qquad \text{（\eqref{eq:dualsym}）}
\end{equation}

\subsection{英汉术语表}

以下按主题分组。请在听课前把\textbf{英文}一列读熟——
课堂上老师大概率直接用英文说这些词。

\begin{table}[htbp]
\centering
\caption{英汉术语表（一）：对称性与算符。}
\label{tab:glossary1}
\begin{tabular}{@{}lll@{}}
\toprule
英文 & 中文 & 本讲义位置 \\
\midrule
generalized global symmetry & 广义全局对称性 & 第 \ref{sec:pform} 节 \\
$q$-form symmetry & $q$-form 对称性／$q$ 阶形式对称性
  & 定义 \ref{def:qform} \\
higher-form symmetry & 高阶形式对称性（$q\ge1$） & 第 \ref{sec:pform} 节 \\
higher symmetry & 高对称性（统称） & 第 \ref{subsec:beyond} 小节 \\
global symmetry & 全局对称性 & 第 \ref{subsec:global-vs-gauge} 小节 \\
gauge symmetry / redundancy & 规范对称性／描述冗余 & 注意 \ref{warn:gauge} \\
continuous / discrete symmetry & 连续／离散对称性 & 第 \ref{sec:ordinary} 节 \\
symmetry operator & 对称算符 & 要点 \ref{kp:U-operator} \\
charged operator & 带荷算符 & 第 \ref{sec:charges-defects} 节 \\
local operator & 局域算符 & 第 \ref{sec:operators} 节 \\
extended operator & 拓展算符 & 第 \ref{subsec:extended} 小节 \\
line / surface operator & 线／面算符 & 第 \ref{subsec:extended} 小节 \\
genuine operator & 真算符（不需依附高维算符） & 第 \ref{subsec:genuine} 小节 \\
Wilson line / loop & Wilson 线／Wilson 圈 & 例 \ref{ex:wilson} \\
't Hooft line & 't Hooft 线 & 例 \ref{ex:thooft} \\
topological operator & 拓扑算符 & 定义 \ref{def:topological} \\
topological defect line (TDL) & 拓扑缺陷线 & 第 \ref{subsec:isingTDL} 小节 \\
defect / twist defect & 缺陷／扭曲缺陷 & 第 \ref{subsec:defects} 小节 \\
domain wall & 畴壁 & 例 \ref{ex:domainwall} \\
twisted boundary condition & 扭曲边界条件 & \eqref{eq:twisted} \\
twisted sector & 扭曲扇区 & 第 \ref{subsec:defects} 小节 \\
fusion rule / fusion category & 融合规则／融合范畴 & \eqref{eq:fusion} \\
invertible / non-invertible & 可逆／不可逆 & 要点 \ref{kp:noninvertible-first} \\
quantum dimension & 量子维数 & 要点 \ref{kp:noninvertible-first} \\
center symmetry & 中心对称性 & 第 \ref{subsec:center} 小节 \\
$N$-ality & $N$ 度 & 例 \ref{ex:Nality} \\
screening & 屏蔽 & 例 \ref{ex:screening} \\
gauging & 规范化（把对称性变成规范对称性） & 第 \ref{subsec:gauging} 小节 \\
global structure & 全局结构 & 要点 \ref{kp:globalstructure} \\
\bottomrule
\end{tabular}
\end{table}

\begin{table}[htbp]
\centering
\caption{英汉术语表（二）：几何、拓扑与反常。}
\label{tab:glossary2}
\begin{tabular}{@{}lll@{}}
\toprule
英文 & 中文 & 本讲义位置 \\
\midrule
differential form / $p$-form & 微分形式／$p$-形式
  & 第 \ref{subsec:forms} 小节 \\
wedge product & 楔积 & \eqref{eq:wedge} \\
exterior derivative & 外微分 & \eqref{eq:dsquared} \\
closed / exact form & 闭形式／恰当形式 & 第 \ref{subsec:forms} 小节 \\
Hodge star & Hodge 星号 & \eqref{eq:hodge} \\
Stokes' theorem & Stokes 定理 & \eqref{eq:stokes} \\
de Rham cohomology & de Rham 上同调 & \eqref{eq:deRham} \\
homologous / homology class & 同调／同调类 & \eqref{eq:homologyinv} \\
codimension & 余维 & 定义 \ref{def:qform} \\
linking number & 连接数 & 定义 \ref{def:link} \\
Poincaré dual & Poincaré 对偶 & \eqref{eq:poincare} \\
Bianchi identity & Bianchi 恒等式 & \eqref{eq:bianchi} \\
Dirac quantization & Dirac 量子化 & \eqref{eq:diracphase} \\
flux & 通量 & 注 \ref{rmk:gauss} \\
Aharonov--Bohm phase & Aharonov--Bohm 相位 & \eqref{eq:diracphase} \\
Ward identity & Ward 恒等式 & \eqref{eq:ward} \\
order parameter & 序参量 & 要点 \ref{kp:ssb-0form} \\
spontaneous symmetry breaking (SSB) & 自发对称破缺
  & 第 \ref{subsec:ssb-unified} 小节 \\
Goldstone boson & Goldstone 玻色子 & 要点 \ref{kp:goldstone-general} \\
area law / perimeter law & 面积律／周长律 & \eqref{eq:arealaw} \\
confinement / deconfinement & 禁闭／解禁闭 & 要点 \ref{kp:confinement} \\
't Hooft anomaly & 't Hooft 反常 & 定义 \ref{def:thooft} \\
mixed anomaly & 混合反常 & 要点 \ref{kp:mutual} \\
anomaly inflow & 反常流入 & 第 \ref{subsec:thooft} 小节 \\
anomaly matching & 反常匹配 & 定理 \ref{thm:matching} \\
background (gauge) field & 背景（规范）场 & 第 \ref{subsec:thooft} 小节 \\
local counterterm & 局域反项 & 定义 \ref{def:thooft} \\
SPT phase & 对称性保护拓扑相 & 第 \ref{subsec:thooft} 小节 \\
electromagnetic duality & 电磁对偶 & \eqref{eq:EMduality} \\
Kramers--Wannier duality & Kramers--Wannier 对偶 & 要点 \ref{kp:KW} \\
T-duality & T 对偶 & \eqref{eq:Tduality} \\
winding number & 缠绕数 & 第 \ref{subsec:compactscalar} 小节 \\
higher group & 高群 & 第 \ref{subsec:beyond} 小节 \\
\bottomrule
\end{tabular}
\end{table}

% =====================================================================
\section{常见困惑}
\label{sec:confusions}

以下是初学广义对称性时最常见的十四个困惑，
每条都给出\textbf{可操作的}回答（而不只是“多想想就懂了”）。

\begin{enumerate}[label=\textbf{Q\arabic*.},ref=\arabic*,leftmargin=1.3cm]

% -------------------------------------------------------------
\item \textbf{“对称性 $=$ 拓扑算符”不是循环定义吗？}
\label{q:circular}

不是。这句话的逻辑方向是：
\emph{原来}我们用“群作用在场上使作用量不变”定义对称性，
然后\emph{推导出}对称算符是拓扑的（定理 \ref{thm:topological}）。
\emph{现在}我们把这个推论\textbf{提升为定义}。

之所以可以这样做，是因为“拓扑算符”这个概念\textbf{独立可判定}：
给你一个算符，你可以直接检查
“它的关联函数在支撑集形变下是否不变”（定义 \ref{def:topological}）。
这个检查不需要事先知道有什么群。

\textbf{收益}：新定义严格更宽——
它自动包含离散对称性（无流）、$q\ge1$ 的情形（非余维 1）、
以及不可逆情形（无群结构），而旧定义都装不下。

% -------------------------------------------------------------
\item \textbf{规范对称性到底算不算对称性？}
\label{q:gauge}

在本课程语境下：\textbf{不算}。见注意 \ref{warn:gauge}。
判别方法：问“它是否作用在物理 Hilbert 空间上非平凡地？”
\begin{itemize}
  \item 全局对称性：$\Uop|\psi\rangle$ 一般是\emph{不同}的物理态
        （如荷不同的态），$\Uop$ 有非平凡谱。
  \item 规范变换：作用在物理态上是\textbf{恒等}的
        （物理态本来就规范不变），什么也没做。
\end{itemize}
不过规范对称性极其重要，因为它\textbf{制造}拓展算符与
高阶形式对称性（要点 \ref{kp:gauge-produces}）。
一句口号：\emph{规范冗余不是对称性，但它是广义对称性的产房。}

% -------------------------------------------------------------
\item \textbf{对称算符、带荷算符、缺陷，我总是搞混。}
\label{q:three}

用\textbf{三个问题}区分（见第 \ref{sec:charges-defects} 节的表）：
\begin{enumerate}[label=(\roman*)]
  \item \emph{它是拓扑的吗？} 是 $\Rightarrow$ 对称算符或缺陷；
        否 $\Rightarrow$ 带荷算符。
  \item \emph{它的支撑集闭合吗？} 闭合 $\Rightarrow$ 对称算符；
        带边 $\Rightarrow$ 缺陷。
  \item \emph{它的维数是多少？} $d-q-1$ $\Rightarrow$ 对称算符侧；
        $q$ $\Rightarrow$ 带荷侧。
\end{enumerate}
在 Maxwell 例子里：$\oint_{\Ssurf_2}\hs F$ 是拓扑的、$\Ssurf_2$ 闭合、
维数 2 $=d-q-1$ $\Rightarrow$ 对称算符；
$W_n(C)$ 不是拓扑的（有面积／周长律）、维数 1 $=q$ $\Rightarrow$ 带荷算符。

% -------------------------------------------------------------
\item \textbf{为什么高阶形式对称性一定阿贝尔？}
\label{q:abelian}

见命题 \ref{prop:abelian}。核心直觉是\textbf{余维}：
\begin{itemize}
  \item 余维 1（$q=0$）：两张“墙”把时空切成有先后的层，
        乘法有顺序 $\Rightarrow$ 可以非阿贝尔。
  \item 余维 $\ge2$（$q\ge1$）：两个支撑集可以在多余的方向上
        \textbf{绕开彼此}，因此顺序无意义 $\Rightarrow$ 必阿贝尔。
\end{itemize}
自查：$d=3$ 中两条直线（余维 2）能不能不相交地交换位置？能。
$d=3$ 中两个平面（余维 1）能不能？不能，它们必相交。$\checkmark$

% -------------------------------------------------------------
\item \textbf{不可逆的东西凭什么叫“对称性”？}
\label{q:noninv}

因为对称性的\textbf{用处}不需要可逆性。
我们用对称性做三件事：
（i）给关联函数选择定则；（ii）组织谱／Hilbert 空间；
（iii）通过反常约束低能物理。
这三件事只需要\textbf{“存在可自由移动的拓扑算符”}。

具体看 Ising 的 $\mathcal N$（第 \ref{subsec:isingTDL} 小节）：
它在 $\sigma$ 上的本征值是 $0$，
于是把 $\mathcal N$ 收缩到包围 $\sigma$ 的圈上就得到
“某些关联函数必须为零”——这是一个\textbf{真正的选择定则}，
和 $\ZZ_2$ 给出的选择定则同样有用，只不过 $\mathcal N$ 没有逆。

\textbf{历史类比}：当年把“反演 $P$”“时间反演 $T$”这些离散变换
也叫对称性时，同样有人反对（“它们没有守恒流”）。
概念的边界随着用处而扩张，这是正常的。

% -------------------------------------------------------------
\item \textbf{1-form 对称性自发破缺，Goldstone 模在哪？}
\label{q:goldstone}

在\textbf{光子}里。见要点 \ref{kp:goldstone-general}：
$q$-form 对称性破缺给出无质量的 $q$-form 规范场。
$q=1$ 时那就是一个无质量的 1-form 场 $A_\mu$，即光子。

初学者常期待“又一个无质量标量”，
但那只对 $q=0$ 成立。破缺的是 1-form 对称性，
Goldstone 模的“形式次数”也随之升高。
数一下自由度也自洽：4 维无质量光子 2 个横向极化，
对应两个 1-form 对称性（电、磁）各自破缺。

% -------------------------------------------------------------
\item \textbf{$(d-1)$-form 对称性是什么？算符成了点？}
\label{q:dminus1}

对，$q=d-1$ 时对称算符维数是 $d-(d-1)-1=0$，即支撑在\textbf{时空点}上。
这意味着存在\textbf{局域的拓扑算符}。

这种情形很特殊：局域拓扑算符 $\Uop(x)$ 因拓扑性而与位置无关，
可以同时对角化，其本征值\textbf{标记了互不通信的“宇宙”}
（\en{universes}）——理论\textbf{分解}（\en{decomposition}）成若干
超选择扇区。带荷对象是 $(d-1)$ 维的，即畴壁。

物理例子：2 维 $\Uone$ 规范理论、
或存在拓扑 $\theta$ 项且真空按 $\theta$ 标记的情形。
第一次听课不必深究，但要认得这个词。

% -------------------------------------------------------------
\item \textbf{为什么我学 QFT 时从没遇到 1-form 对称性？}
\label{q:whynever}

见要点 \ref{kp:whypure}。三个原因叠加：
\begin{enumerate}[label=(\roman*)]
  \item 标准模型有电荷 $1$ 的物质（电子、夸克），
        把电 1-form 对称性完全破坏（\eqref{eq:UtoZq} 取 $q=1$）；
  \item 没有观测到磁单极子，磁 1-form 对称性虽然存在但用不上；
  \item 传统教科书只讲“作用在粒子态上的对称性”，
        而 1-form 对称性作用在线算符上，
        在散射振幅语言里看不见。
\end{enumerate}
所以这不是你的知识漏洞，而是\textbf{历史与侧重点}的问题。

% -------------------------------------------------------------
\item \textbf{有些文献说 higher-form symmetry 是“涌现的／近似的”，
什么意思？}
\label{q:emergent}

意思是它只在某个能量窗口或某个近似下成立。例如：
\begin{itemize}
  \item 若电荷 $q$ 很大，$\ZZ_q^{(1)}$ 虽是精确的，
        但在低能有效理论中常可近似当作 $\Uone^{(1)}$；
  \item 磁 1-form 对称性在有磁单极子的理论中被破坏，
        但若单极子很重，在低能下\textbf{近似}守恒
        （磁流守恒在磁流体力学中就是这样用的）。
\end{itemize}
判断方法照旧：\emph{对称算符在多好的近似下是拓扑的？}
$\ext\hs J$ 有多小？

% -------------------------------------------------------------
\item \textbf{Wilson 线是对称算符还是带荷算符？}
\label{q:wilsonrole}

在纯 Maxwell 中它是\textbf{带荷算符}（被 $\Uop_e$ 作用）。
这是最常见的混淆点，因为 $W_n=e^{in\oint A}$ 的形式
与对称算符 $e^{i\alpha\oint \hs F/e^2}$ 长得很像。
区别是被积的东西：
\begin{center}
$A$（不闭：$\ext A=F\ne0$）$\Rightarrow$ 不拓扑 $\Rightarrow$ 带荷；\qquad
$\hs F$（闭：$\ext\hs F=0$）$\Rightarrow$ 拓扑 $\Rightarrow$ 对称算符。
\end{center}
\textbf{“看被积形式闭不闭”是最快的判别法。}

% -------------------------------------------------------------
\item \textbf{“$p$-form 对称性”里的 $p$ 到底数什么？}
\label{q:whichp}

数\textbf{带荷对象的维数}（本讲义的 $q$）。
这是 GKSW 的标准约定，也是绝大多数文献的约定。
但\textbf{偶有作者用对称算符的余维或电流次数来命名}，
读文献时务必用要点 \ref{kp:dimensioncheck} 的三连自查
反推作者的约定，不要硬记名字。

% -------------------------------------------------------------
\item \textbf{'t Hooft 反常和我学过的 ABJ 反常是一回事吗？}
\label{q:anomalykinds}

不是，见注意 \ref{warn:twoanomalies}。
\begin{itemize}
  \item ABJ：对称性\textbf{被破坏}，$\partial_\mu j^\mu\ne0$，
        守恒律没了。
  \item 't Hooft：对称性\textbf{完好}，$\partial_\mu j^\mu=0$；
        只是\textbf{不能规范化}。
\end{itemize}
一个统一的看法：ABJ 反常可以理解为“把某个对称性规范化之后，
另一个对称性出现了 't Hooft 反常”，
所以两者在背景场语言中是同一个框架的不同切片。
本课程主要用第二种。

% -------------------------------------------------------------
\item \textbf{规范化（gauging）与自发破缺有什么关系？}
\label{q:gauge-vs-ssb}

两者\textbf{完全不同}，但常被并列提到：
\begin{center}
\begin{tabular}{@{}lll@{}}
\toprule
 & 自发破缺 & 规范化 \\
\midrule
是否改变理论 & 否（同一个理论的一个相） & \textbf{是}（得到\emph{新}理论） \\
对称性去哪了 & 仍存在，只是真空不对称 & 变成规范冗余，不再是全局对称性 \\
后果 & Goldstone 模、真空退化 & 算符谱被投影，出现对偶对称性 \\
可否总是做 & 由动力学决定 & 有 't Hooft 反常时\textbf{不可以} \\
\bottomrule
\end{tabular}
\end{center}

% -------------------------------------------------------------
\item \textbf{为什么老师总要把理论放在一般（弯曲、非平凡拓扑）流形上？}
\label{q:manifold}

因为\textbf{广义对称性的信息藏在拓扑里}。
在 $\RR^d$ 上，所有闭曲面都可以收缩到一点，
$H^p(\RR^d)$ 除 $p=0$ 外全为零，
于是所有连接数都是零、所有对称算符都平凡——\textbf{什么都看不见}。

只有把时空取成有非平凡拓扑的（如 $S^1\times S^3$、$T^4$、
或挖去一条线的 $\RR^4$），
才能有非零的连接数、非平凡的通量扇区、
以及区分不同全局结构的观测量。
所以“放到一般流形上”不是数学炫技，而是\textbf{测量工具}。

\end{enumerate}

% =====================================================================
\section{听课重点}
\label{sec:focus}

\subsection{三条必须听懂的主线}

\begin{enumerate}[label=\textbf{\Roman*.},leftmargin=1.1cm]
  \item \textbf{“对称性 $=$ 拓扑算符”的完整逻辑。}
        听老师如何从 $\partial_\mu j^\mu=0$ 走到
        “$\Uop(\Ssurf)$ 只依赖 $[\Ssurf]$”。
        若这一步没跟上，后面全部内容都会变成术语堆砌。
        \textbf{预习对策}：第 \ref{subsec:topological-rewrite} 小节。
  \item \textbf{维数记账。}
        每当老师写下一个对称算符，立刻在笔记边上标注
        “$d=?$，$q=?$，支撑维数 $=?$，带荷对象维数 $=?$”。
        \textbf{预习对策}：要点 \ref{kp:dimensioncheck}、表 \ref{tab:pform}。
  \item \textbf{Maxwell 的两个 1-form 对称性。}
        这是全课程的标准例子，后面的一切
        （反常、对偶、全局结构、禁闭）都会回到它。
        \textbf{预习对策}：第 \ref{subsec:maxwell} 小节，
        务必自己推一遍 \eqref{eq:fluxisn}。
\end{enumerate}

\subsection{课上要主动记录的六件事}

\begin{enumerate}[label=\textbf{\arabic*.},leftmargin=1.1cm]
  \item \textbf{老师的约定}。$q$ 数什么？$2\pi$ 放哪？
        号差与信号（洛伦兹／欧氏）？
        \textbf{第一堂课就把约定抄下来}，
        否则后面对不上因子会很痛苦。
  \item \textbf{每个例子的“三件套”}：
        对称算符是什么、带荷对象是什么、荷怎么读出来。
        听完一个例子若答不出这三问，就是没听懂。
  \item \textbf{哪些结论依赖运动方程，哪些是恒等式}
        （要点 \ref{kp:two-sources}）。
        这决定了加入物质场后对称性是否存活。
  \item \textbf{每次出现“不能规范化”时，反常是什么形式}。
        记下 $(d+1)$ 维拓扑项的形状，哪怕系数记不准。
  \item \textbf{融合规则}。看到 $\mathcal N\times\mathcal N=\cdots$
        这类式子立刻判断可逆性，并用“本征值相乘”自查
        （见第 \ref{subsec:isingTDL} 小节练习）。
  \item \textbf{老师提到的开放问题}。这个领域仍在快速发展，
        老师随口说的“这个还不清楚”往往是很好的课题线索。
\end{enumerate}

\subsection{自检清单：上课前能否回答}

\begin{enumerate}[label=$\square$\ ,leftmargin=1.1cm]
  \item 我能在 5 行内推出 Noether 流 \eqref{eq:noethercurrent}。
  \item 我能说明为什么 $\ext\hs j=0$ 等价于
        $\Uop(\Ssurf)$ 可自由形变。
  \item 给定 $d$ 与 $q$，我能立刻写出对称算符与带荷对象的维数，
        并验证 $(d-q-1)+q+1=d$。
  \item 我能写出 Maxwell 的两个对称算符 \eqref{eq:maxwellops}
        并说清各自的守恒来源。
  \item 我能推导 $\frac{1}{e^2}\oint_{\Ssurf_2}\hs F=n$
        （即 \eqref{eq:fluxisn}）。
  \item 我知道为什么加入电荷 $q$ 的物质把
        $\Uone^{(1)}_e$ 打成 $\ZZ_q^{(1)}$。
  \item 我能陈述“禁闭 $\Leftrightarrow$ 1-form 对称性未破缺”
        并说出对应的 Wilson 圈行为。
  \item 我能区分对称算符、带荷算符、缺陷（Q\ref{q:three} 的三个问题）。
  \item 我能从 $\mathcal N^2=\mathbf 1+\eta$ 论证 $\mathcal N$ 不可逆。
  \item 我知道 't Hooft 反常与 ABJ 反常的区别。
  \item 我能说出规范化 $q$-form $G$ 后得到什么对称性
        （\eqref{eq:dualsym}）。
  \item 我理解为什么要在非平凡拓扑的流形上讨论（Q\ref{q:manifold}）。
\end{enumerate}

\subsection{课后 48 小时内应该做的事}

\begin{itemize}
  \item 把课上每个新对称性填进一张自制的表 \ref{tab:pform} 式字典里；
  \item 至少自己重做一遍课上的连接数／通量计算；
  \item 把老师用的约定与本讲义的约定做一次对照表，
        标出差异（这一步能省掉后面大量困惑）。
\end{itemize}

% =====================================================================
\section{进一步学习路线}
\label{sec:further}

\subsection{建议顺序}

\begin{enumerate}[label=\textbf{阶段 \arabic*：},leftmargin=1.9cm]
  \item \textbf{奠基（读原始论文）}。
        Gaiotto--Kapustin--Seiberg--Willett，
        \emph{Generalized Global Symmetries}（arXiv:1412.5148）。
        这是本领域的\textbf{奠基文献}，通常简称 \textbf{GKSW}。
        建议读第 1--3 节：那里就是本讲义第
        \ref{sec:pform}--\ref{subsec:maxwell} 小节的内容，
        只是更简洁。读完你会发现本讲义的 Maxwell 推导
        与其第 2 节几乎逐行对应。
  \item \textbf{本课程直接相关的讲义}。
        D.~Brennan 与 S.~Hong，
        \emph{Introduction to Generalized Global Symmetries in QFT
        and Particle Physics}（arXiv:2306.00912）。
        \textbf{授课老师本人合写的入门讲义}，
        侧重与粒子物理的联系。若只读一份补充材料，读这份。
  \item \textbf{全景综述}。
        \begin{itemize}
          \item Córdova--Dumitrescu--Intriligator--Shao，
                Snowmass 白皮书
                \emph{Generalized Symmetries in Quantum Field Theory
                and Beyond}（arXiv:2205.09545）——最好的“地图”。
          \item J.~McGreevy，
                \emph{Generalized Symmetries in Condensed Matter}
                （arXiv:2204.03045）——凝聚态视角，
                对格点模型与 Ising 例子讲得很实在。
        \end{itemize}
  \item \textbf{不可逆／范畴型对称性}。
        \begin{itemize}
          \item S.~Schäfer-Nameki，
                \emph{ICTP Lectures on (Non-)Invertible Generalized
                Symmetries}（arXiv:2305.18296）。
          \item S.-H.~Shao，
                \emph{What's Done Cannot Be Undone: TASI Lectures on
                Non-Invertible Symmetries}（arXiv:2308.00747）——
                Ising 的 $\mathcal N$ 线在这里讲得最透。
        \end{itemize}
  \item \textbf{反常与应用}。
        \begin{itemize}
          \item Gaiotto--Kapustin--Komargodski--Seiberg，
                \emph{Theta, Time Reversal, and Temperature}
                （arXiv:1703.00501）——
                $\ZZ_N$ 1-form 中心对称性与 $\theta$ 角混合反常的
                标志性应用。
          \item Kapustin--Seiberg，
                \emph{Coupling a QFT to a TQFT and Duality}
                （arXiv:1401.0740）——GKSW 的前身，
                讲清了“全局结构”的意思。
          \item 't Hooft 1980 年的经典文章
                \emph{Naturalness, Chiral Symmetry, and Spontaneous
                Chiral Symmetry Breaking}——反常匹配论证的源头。
        \end{itemize}
\end{enumerate}

\begin{warning}
上面的 arXiv 编号供检索定位之用。文献版本、题名与编号偶有变动，
\textbf{请以老师课上给出的书目为准}，
并在 arXiv 上按作者与题名再确认一次。
\end{warning}

\subsection{按主题的补充线索}

\begin{center}
\begin{tabular}{@{}ll@{}}
\toprule
想深入的主题 & 关键词（用于检索） \\
\midrule
高群 & higher group, 2-group symmetry, Córdova--Dumitrescu--Intriligator \\
对称拓扑场论 & symmetry topological field theory, SymTFT, sandwich construction \\
不可逆手征对称性 & non-invertible chiral symmetry, axion, Choi--Córdova--Hsin--Lam--Shao \\
不可逆对偶缺陷 & duality defect, Tambara--Yamagami, Kaidi--Ohmori--Zheng \\
2 维 CFT 拓扑线 & Verlinde line, topological defect line, Petkova--Zuber \\
格点起源 & Wegner $\ZZ_2$ gauge theory, Wilson loop order parameter \\
凝聚态应用 & fracton, subsystem symmetry, toric code, Nussinov--Ortiz \\
反常的数学 & Freed--Moore--Segal, invertible field theory, cobordism \\
超对称理论中的应用 & $\mathcal N=4$ global structure, $\SUgrp(N)$ vs $\mathrm{PSU}(N)$ \\
\bottomrule
\end{tabular}
\end{center}

\subsection{三个动手项目（巩固效果最好）}

\begin{enumerate}[label=\textbf{项目 \arabic*.},leftmargin=1.9cm]
  \item \textbf{格点 $\ZZ_2$ 规范理论的 1-form 对称性}。
        在 3 维（或 4 维）立方格点上写下 Wegner 的 $\ZZ_2$ 规范理论，
        显式构造 1-form 对称算符（是链上 $\sigma^x$ 的乘积，
        支撑在对偶格点的闭曲面上），
        验证它与 Hamilton 量交换，
        并计算 Wilson 圈在强／弱耦合极限下的面积／周长律。
        \textbf{这是把本讲义所有抽象概念变成有限维线性代数的最好练习。}
  \item \textbf{Ising 链上显式构造 KW 对偶算符}。
        用 \eqref{eq:KWvars} 在有限链上写出对偶变换矩阵，
        数值验证临界点 $J=h$ 处的谱自对偶，
        并检查 $\mathcal N$ 在 $\ZZ_2$ 奇态上的“零本征值”。
  \item \textbf{Maxwell 通量扇区}。
        把 4 维 Maxwell 放在 $T^2\times\RR^2$ 上，
        计算磁通 $\frac{1}{2\pi}\oint_{T^2}F=m$ 的各个扇区，
        显式看出磁 1-form 对称性如何标记这些扇区。
\end{enumerate}

% =====================================================================
\appendix

\section{附录：微分形式与拓扑速查}
\label{app:forms}

供做题时快速回查。

\subsection{次数记账表}

\begin{center}
\begin{tabular}{@{}lll@{}}
\toprule
对象 & 形式次数 & 可积于 \\
\midrule
函数 $f$ & 0 & 点 \\
规范场 $A$ & 1 & 曲线 $C$ \\
场强 $F=\ext A$ & 2 & 曲面 $\Ssurf_2$ \\
$\hs F$（$d=4$） & 2 & 曲面 $\Ssurf_2$ \\
$q$-form 对称性的电流 $J$ & $q+1$ & --- \\
$\hs J$ & $d-q-1$ & $\Mfd_{d-q-1}$ \\
$k$ 维子流形的 $\delta$ 形式 $\delta_{\Mfd_k}$ & $d-k$ & --- \\
体积元 $\Vol_d$ & $d$ & 整个时空 \\
\bottomrule
\end{tabular}
\end{center}

\subsection{常用恒等式}

\begin{align}
  \ext^2&=0 \\
  \ext(\omega_{(p)}\wedge\eta_{(q)})
    &=\ext\omega_{(p)}\wedge\eta_{(q)}
     +(-1)^p\omega_{(p)}\wedge \ext\eta_{(q)} \\
  \omega_{(p)}\wedge\eta_{(q)}&=(-1)^{pq}\eta_{(q)}\wedge\omega_{(p)} \\
  \int_{\Mfd_{p+1}}\ext\omega_{(p)}
    &=\int_{\partial\Mfd_{p+1}}\omega_{(p)}
    &&\text{(Stokes)}\\
  \hs\hs\,\omega_{(p)}&=(-1)^{p(d-p)}\omega_{(p)}
    &&\text{(欧氏号差)}\\
  \partial_{\mu_1}J^{\mu_1\cdots\mu_{q+1}}=0
    &\ \Longleftrightarrow\ \ext\hs J=0 \\
  \int_{\Mfd_d}\delta_{\Mfd_k}\wedge\omega_{(k)}
    &=\int_{\Mfd_k}\omega_{(k)}
    &&\text{(Poincaré 对偶)}\\
  \int_{V_{k}}\delta_{\Mfd_{d-k}}
    &=\#\big(V_k\cap \Mfd_{d-k}\big)
    &&\text{(相交数)}
\end{align}

\subsection{维数自查流程}

给定一个陈述“$\Uop$ 是 $q$-form 对称算符，作用在 $\Op$ 上”：

\begin{enumerate}[label=\textbf{步骤 \arabic*.},leftmargin=1.9cm]
  \item 确认时空维数 $d$。
  \item 计算 $\dim\Uop=d-q-1$，与题中给的支撑维数核对。
  \item 计算 $\dim\Op=q$，与题中给的带荷对象核对。
  \item 检查 $\dim\Uop+\dim\Op+1=d$。
  \item 若用连续对称性，检查电流是 $(q+1)$-形式、$\hs J$ 是
        $(d-q-1)$-形式。
\end{enumerate}

任一步不通过，就是记错了约定或记错了维数——回到
要点 \ref{kp:dimensioncheck} 重新数。

\subsection{三个例子的完整参数表}

\begin{center}
\begin{tabular}{@{}llll@{}}
\toprule
 & Ising（$\ZZ_2$） & Maxwell（电） & Maxwell（磁）\\
\midrule
$d$ & 2 & 4 & 4 \\
$q$ & 0 & 1 & 1 \\
群 & $\ZZ_2$ & $\Uone$ & $\Uone$ \\
$\hs J$ & --- (离散) & $\hs F/e^2$ & $F/2\pi$ \\
守恒来源 & --- & 运动方程 & Bianchi \\
对称算符维数 & 1 & 2 & 2 \\
对称算符 & $\prod_j\sigma^x_j$
  & $e^{i\alpha\oint\hs F/e^2}$ & $e^{i\beta\oint F/2\pi}$ \\
带荷对象维数 & 0 & 1 & 1 \\
带荷对象 & $\sigma^z_j$ & $W_n(C)$ & $T_m(C)$ \\
荷 & $\pm1$ & $n\in\ZZ$ & $m\in\ZZ$ \\
破缺相 & 有序相 & Coulomb 相 & Coulomb 相 \\
未破缺相 & 无序相 & 禁闭相 & --- \\
\bottomrule
\end{tabular}
\end{center}

\subsection{练习答案要点}

\begin{itemize}
  \item \textbf{练习 \ref{ex:forms}(a)}：
        由 \eqref{eq:dsquared} 取 $p=1$，
        $\ext(A\wedge F)=\ext A\wedge F-A\wedge \ext F
        =F\wedge F-0=F\wedge F$。
        故 $\int F\wedge F$ 是全导数的积分，只依赖拓扑。
  \item \textbf{练习 \ref{ex:pivot}(b)}：
        $Q(\Ssurf)-Q(\Ssurf')=q$，即扫过一个荷 $q$ 的算符时
        荷读数跳变 $q$；等价地
        $\Uop_\alpha(\Ssurf)=e^{iq\alpha}\Uop_\alpha(\Ssurf')$
        （作用在含该算符的关联函数上）。
  \item \textbf{练习 \ref{ex:gaugecheck}}：
        (a) 2-form $\ZZ_N$；(b) 1-form $\ZZ_N$（形式不变）；
        (c) 3 维中 0-form $\ZZ_N$ 规范化得 1-form $\ZZ_N$，
        两者的带荷对象（局域算符与线算符）之间的相互相位
        $e^{2\pi i nm/N}$ 就是 $\ZZ_N$ 规范理论中
        电／磁线的辫结相位。
  \item \textbf{练习 \ref{ex:maxwell}(c)}：
        $N$ 度为 2 的物质能屏蔽 $n_R\in2\ZZ$ 的线，
        残留 $\ZZ_N/\langle 2\rangle=\ZZ_{\gcd(2,N)}$；
        $N$ 偶时残留 $\ZZ_2$，$N$ 奇时完全破坏。
\end{itemize}

\subsection{一页总结}

\begin{center}
\fbox{\parbox{0.92\textwidth}{
\textbf{广义对称性一页总结}
\begin{enumerate}[leftmargin=1.2cm,itemsep=3pt]
  \item 对称性 $=$ \textbf{拓扑算符}（可自由形变的算符）。
        对连续情形，拓扑性 $\Leftrightarrow$ $\ext\hs J=0$
        $\Leftrightarrow$ 守恒律。
  \item $q$-form 对称性：算符在 $(d-q-1)$ 维闭流形上，
        带荷对象 $q$ 维，$(d-q-1)+q+1=d$。
  \item 荷 $=$ 连接数读出的相位，不是贴在算符上的标签。
  \item 闭合支撑 $\Rightarrow$ 对称算符；带边支撑 $\Rightarrow$ 缺陷，
        其边界是物理激发。
  \item $q\ge1$ 必阿贝尔（余维 $\ge2$ 可绕开）。
  \item Maxwell：运动方程给电 $\Uone^{(1)}_e$（作用在 Wilson 线上），
        Bianchi 给磁 $\Uone^{(1)}_m$（作用在 't Hooft 线上）。
  \item 加入电荷 $q$ 物质：$\Uone^{(1)}_e\to\ZZ^{(1)}_q$；
        Bianchi 不受影响。
  \item 1-form 对称性未破缺（面积律）$=$ 禁闭；
        破缺（周长律）$=$ 解禁闭，Goldstone 模就是光子。
  \item 融合无逆元 $\Rightarrow$ 不可逆对称性
        （Ising：$\mathcal N^2=\mathbf 1+\eta$，$\mathcal N\sigma=0$）。
  \item 't Hooft 反常 $=$ 不能规范化；它是 RG 不变量，
        因此约束低能物理（破缺／无能隙／TQFT 三选一）。
\end{enumerate}
}}
\end{center}

\end{document}
